
\documentclass[10pt, reqno]{amsart}
\usepackage{amsmath, amsthm, amscd, amsfonts, amssymb, graphicx, color}
\usepackage[bookmarksnumbered, colorlinks, plainpages]{hyperref}
\usepackage{quiver}
\textheight 22.5truecm \textwidth 14.5truecm
\setlength{\oddsidemargin}{0.35in}\setlength{\evensidemargin}{0.35in}
\definecolor{MK_One_Five}{RGB}{90,180,172}
\definecolor{MK_Two_Six}{RGB}{33,102,172}
\definecolor{MK_Three_Six}{RGB}{27,120,55}
\definecolor{MK_Three_One}{RGB}{118,42,131}
\hypersetup{linkcolor=MK_Three_Six,citecolor=MK_Two_Six,urlcolor=MK_Three_One}
\setlength{\topmargin}{-.5cm}

\makeatletter
%Table of Contents
\setcounter{tocdepth}{2}

% Add bold to \section titles in ToC and remove . after numbers
\renewcommand{\tocsection}[3]{%
  \indentlabel{\@ifnotempty{#2}{\bfseries\ignorespaces#1 #2\quad}}\bfseries#3}
% Remove . after numbers in \subsection
\renewcommand{\tocsubsection}[3]{%
  \indentlabel{\@ifnotempty{#2}{\ignorespaces#1 #2\quad}}#3}
%\let\tocsubsubsection\tocsubsection% Update for \subsubsection
%...

\newcommand\@dotsep{4.5}
\def\@tocline#1#2#3#4#5#6#7{\relax
  \ifnum #1>\c@tocdepth % then omit
  \else
    \par \addpenalty\@secpenalty\addvspace{#2}%
    \begingroup \hyphenpenalty\@M
    \@ifempty{#4}{%
      \@tempdima\csname r@tocindent\number#1\endcsname\relax
    }{%
      \@tempdima#4\relax
    }%
    \parindent\z@ \leftskip#3\relax \advance\leftskip\@tempdima\relax
    \rightskip\@pnumwidth plus1em \parfillskip-\@pnumwidth
    #5\leavevmode\hskip-\@tempdima{#6}\nobreak
    \leaders\hbox{$\m@th\mkern \@dotsep mu\hbox{.}\mkern \@dotsep mu$}\hfill
    \nobreak
    \hbox to\@pnumwidth{\@tocpagenum{\ifnum#1=1\bfseries\fi#7}}\par% <-- \bfseries for \section page
    \nobreak
    \endgroup
  \fi}
\AtBeginDocument{%
\expandafter\renewcommand\csname r@tocindent0\endcsname{0pt}
}
\def\l@subsection{\@tocline{2}{0pt}{2.5pc}{5pc}{}}
\makeatother

\usepackage{lipsum}
\usepackage{parskip}
\usetikzlibrary{shapes.geometric, arrows}

\newtheorem{theorem}{Theorem}[section]
\newtheorem{definition-theorem}[theorem]{Definition-Theorem}
\newtheorem{lemma}[theorem]{Lemma}
\newtheorem{proposition}[theorem]{Proposition}
\newtheorem{prop-def}[theorem]{Proposition-Definition}
\newtheorem{corollary}[theorem]{Corollary}
\newtheorem{conjecture}[theorem]{Conjecture}
\newtheorem{fact}[theorem]{Fact}
\newtheorem{question}[theorem]{Question}

\theoremstyle{definition}
\newtheorem{definition}[theorem]{Definition}
\newtheorem{notation}[theorem]{Notation}
\newtheorem{example}[theorem]{Example}
\newtheorem{exercise}[theorem]{Exercise}
\newtheorem{conclusion}[theorem]{Conclusion}
\newtheorem{criterion}[theorem]{Criterion}
\newtheorem{summary}[theorem]{Summary}
\newtheorem{axiom}[theorem]{Axiom}
\newtheorem{problem}[theorem]{Problem}

\theoremstyle{remark}
\newtheorem{claim}[theorem]{Claim}
\newtheorem{claim*}[theorem]{Claim}
\newtheorem{remark}[theorem]{Remark}
\newtheorem{assumption}[theorem]{Assumption}
\numberwithin{equation}{section}

\newcommand{\ZZ}{\mathbb{Z}}
\newcommand{\QQ}{\mathbb{Q}}
\newcommand{\RR}{\mathbb{R}}
\newcommand{\CC}{\mathbb{C}}
\newcommand{\HH}{\mathbb{H}}
\newcommand{\Gr}{\mathrm{Gr}}
\newcommand{\GL}{\mathrm{GL}}
\newcommand{\GSp}{\mathrm{GSp}}
\newcommand{\M}{\mathrm{M}}%matrix%
\newcommand{\pl}{\mathcal{P}} %poincare bundle%
\newcommand{\pt}{\mathcal{P}^{\times}} %poincare torsor%
\newcommand{\pic}{\mathrm{Pic}}
\newcommand{\gm}{\mathbb{G}_{\mathrm{m}}}
\newcommand{\pr}{\mathrm{pr}}
\newcommand{\id}{\mathrm{id}}
\newcommand{\spec}{\operatorname{Spec}}
\newcommand{\lb}{\mathcal{L}}%line bundle%
\newcommand{\ag}{{S_{g}}} %siegel modular variety%
\newcommand{\Ag}{{\mathcal{A}_g}}%universal abelian scheme%
\newcommand{\Agv}{{\mathcal{A}_g^\vee}}
\newcommand{\Bg}{\mathcal{B}_g}
\newcommand{\Pl}{\mathcal{P}_g} %universal poincare bundle%
\newcommand{\Pt}{\mathcal{P}^{\times}_g} %universal poincare torsor%

\newcommand{\Hom}{\mathrm{Hom}}
\newcommand{\Ext}{\mathrm{Ext}}
\newcommand{\NS}{\mathrm{NS}} %neron severi group
\newcommand{\Ga}{\mathbb {G}_a} %additive group
\newcommand{\Gm}{\mathbb {G}_m} %mult group
\newcommand{\pgn}{\mathcal{P}^{\times}_{g,n}}
\newcommand{\im}{\mathrm {im}}

\newcommand{\Jac}{\mathrm {Jac}}

\newcommand{\calA}{\mathcal A}
\newcommand{\calB}{\mathcal{B}}
\newcommand{\calG}{\mathcal{G}}
\newcommand{\calH}{\mathcal H}
\newcommand{\calO}{\mathcal {O}}
\newcommand{\calM}{\mathcal {M}} %S_alpha,beta(L) construction in trivializing poincare line bundle

\newcommand{\calP}{\mathcal P}

\begin{document}
\setcounter{page}{1}

\noindent 

\centerline{}

\centerline{}

\title[\scalebox{.8}{On Special Subvarieties of the Universal Semi-abelian Scheme and Pink Conjectures}]{On Special Subvarieties of the Universal Semi-abelian Scheme and Pink Conjectures}

\author{Kaiyuan Gu$^1$ and Chenxin Huang$^2$}

\address{$^{1}$ School of Mathematical Sciences, Peking University, Beijing, China.}
\email{\textcolor{MK_One_Five}{gky@stu.pku.edu.cn}}

\address{$^{2}$ School of Mathematical Sciences, Fudan University, Shanghai, China.}
\email{\textcolor{MK_One_Five}{cxhuang21@m.fudan.edu.cn}}

%\dedicatory{This paper is dedicated to Professor ABCD}

\subjclass[2010]{14K05, 14G35, 11G15}

\keywords{Semi-abelian varieties, Poincar\'e biextensions, Mixed Shimura varieties, Pink conjectures, Relative Manin-Mumford conjecture}


\begin{abstract}
The universal Poincar\'{e} torsor, or more generally the universal semi-abelian scheme, can be viewed as a mixed Shimura variety. We give a classification of special subvarieties of the universal semi-abelian scheme of arbitrary toric rank. 

Given this classification, we show that the Zilber-Pink conjecture for mixed Shimura varieties implies the Zilber-Pink conjecture for semi-abelian varieties, correcting an error in an unpublished manuscript of Pink. Moreover, we give a more reasonable reformulation of the Relative Manin-Mumford Conjecture for semi-abelian schemes.

\end{abstract} \maketitle
\parskip=0.5pt
\small
\tableofcontents
\normalsize
\parskip=2pt
\setlength{\parindent}{2em}
\input{section1}
\input{section2}
\input{section3}
\input{section4}
\input{section5}
\input{section6}

\bibliographystyle{amsalpha}
\bibliography{mybibliography}

\end{document}

%------------------------------------------------------------------------------
% End of journal.tex
%------------------------------------------------------------------------------

\section{Introduction}\label{section1}

This paper aims to \textit{classify all special subvarieties in geometric terms} of a given mixed Shimura variety. More precisely, we classified, in Theorem \ref{general}, all special subvarieties of the universal semi-abelian variety with toric part of rank $n$ for all $n\ge 1$ in geometric terms. Then the general case can be deduced by the reduction lemma \cite[Lemma~2.26]{PinkThesis}. 
We achieve our goal mainly relying on the concept of the \textit{Ribet property} (Proposition-Definition \ref{ribet_prop}). Having the classification theorems at hand, we can link \textit{the Relative Manin-Mumford Conjecture} \ref{intro_RMM} with other related conjectures. See Section \ref{section6} for further discussions.

\subsection{Motivations of the Problem}$\ $

We begin with the following motivating question.
\begin{question}
Let $G$ be a commutative algebraic group over an algebraically closed field $k$ of characteristic zero. If $Z\subset G$ is an irreducible subvariety containing a Zariski dense subset of torsion points of $G$, what can one say about the shape of $Z$?
\end{question}
This is a typical example of a \textbf{special point problem}. In the case of abelian varieties, the Manin-Mumford conjecture, proved by Raynaud \cite{Ray83}, claims that such $Z$ is the translate of an abelian subvariety by a torsion point. A similar result holds for semi-abelian varieties.
%There are two famous results of this question when $G$ is the algebraic torus $\gm^n$ or an abelian variety.

%\begin{theorem}[Lang's conjecture, M.Laurent \cite{Lau84}, Sarnak-Adams-Rudnick \cite{SA94}]
%Let $\Sigma$ be a set of torsion points in $\gm^n(\overline{\QQ})$. The Zariski closure of $\Sigma$ is a finite union of translates of algebraic subtori of $\gm^n$ by a torsion point.
%\end{theorem}

%\begin{theorem}[Manin-Mumford conjecture, Raynaud \cite{Ray83}]$ $
%Let $A$ be an abelian variety defined over a field of characteristic 0 and let $\Sigma$ be a set of torsion points in $A$. The Zariski closure of $\Sigma$ is a finite union of translates of abelian subvarieties of $A$ by a torsion point.
%\end{theorem}
In terms of special point problems, for any semi-abelian variety $G$ we call its torsion points the \textit{special points} and the translates of semi-abelian subvarieties by a torsion point are called the \textit{special subvarieties}. Then Raynaud's theorem reads: If $Z\subset G$ is an irreducible subvariety containing a Zariski dense subset of special points of $G$, then $Z$ is a special subvariety.


One can consider the analogous question in the context of Shimura varieties, for instance Siegel moduli varieties, where CM points play the role of special points. The analog of the Manin-Mumford conjecture is the Andr\'{e}-Oort conjecture. There are plenty of works toward this conjecture. The case of Siegel modular varieties is finally proved by Tsimerman \cite{Tsimerman}, based on works of many others, and the general case is claimed in \cite{AO}. 

It is a natural question to attempt to merge these two conjectures. 
An object which unifies both semi-abelian varieties and Shimura varieties is mixed Shimura varieties. Let us introduce the prototype -- the universal Poincar\'{e} torsor, which is also the main object of study in this article. Both the Manin-Mumford conjecture and the Andr\'{e}-Oort conjecture are particular cases of the far-reaching \textit{Zilber-Pink} conjecture for mixed Shimura varietis.

Let $\ag/\mathbb{C}$ be the Siegel modular variety parametrizing principally polarized $g$-dimensional abelian varieties with certain level structures (or more generally of an arbitrary polarization type, see Remark \ref{poltype}). Suppose the level structure is good enough such that $\ag$ is a fine moduli space, whose specific form is unimportant in this article. Then there is a universal abelian scheme $\Ag$ over $\ag$, whose dual is denoted by $\Agv$. One can construct the universal Poincar\'{e} torsor (bi-extension) $\Pt$, which is the $\gm$-torsor coorsponding to the universal Poincar\'{e} line bundle on $\Ag\times_{\ag}\Agv$. We will review the definition in detail in Section \ref{section2}.

The Poincar\'{e} universal torsor is interesting because it has both structures as introduced in the beginning. On one hand, $\Pt$ is a family of semi-abelian varieties over $\Ag$ \textit{and} $\Agv$ (thus the name “bi-extension”). Thus it is a natural object in order to consider a suitable family version of the Manin-Mumford conjecture, called the \textit{relative Manin-Mumford conjecture}. On the other hand, the Siegel modular variety $\ag$ is a typical example of a Shimura variety, and the universal Poincar\'{e} torsor $\Pt$ is a typical example of a \textit{mixed} Shimura variety. One can generalize the notion of CM points in $\ag$ to $\Pt$, and more generally to any mixed Shimura variety, and hence propose the Andr\'{e}-Oort conjecture for mixed Shimura varieties. 




By the work of Gao \cite{GaoAL, GaomixAO}, the Andr\'{e}-Oort conjecture for pure Shimura varieties implies the Andr\'{e}-Oort conjecture for mixed Shimura varieties. The situation becomes much more subtle when it comes to the relative Manin-Mumford conjecture. 

In the unpublished manuscript \cite{Pink05}, Pink proposed a version \cite[Conjecture 6.2]{Pink05} of the relative Manin-Mumford conjecture for semi-abelian schemes. He then used a construction related to $\Pt$ (or more precisely the universal semi-abelian variety) to deduce it from the Zilber-Pink conjecture. However, in \cite{BE20}, the authors constructed a class of counterexamples called the \textit{Ribet sections} to \cite[Conjecture 6.2]{Pink05}. While the authors of \cite{BE20} pointed out the flaws of some arguments in \cite{Pink05}, they also proved that the Ribet sections are special subvarieties of $\Pt$ in the sense of mixed Shimura varieties, thus supporting 
%The counterproof of \cite[Conjecture 6.2]{Pink05} exactly supports
the Zilber-Pink conjecture on unlikely intersections. 
This suggests that when working with semi-abelian schemes, there are more special subvarieties to consider than translates of semi-abelian subschemes. Now the natural question is: should all special subvarieties be constructed from torsion translates of semi-abelian subschemes and Ribet sections?

In this article, we give an affirmative answer to this question. We accomplish this by thoroughly classifying all special subvarieties of $\Pt$, and more generally for the universal semi-abelian variety, viewed as mixed Shimura varieties. For the definition of special subvarieties of mixed Shimura varieties, see \cite[Definition 4.1]{pink04}.  Thus the formulation of relative Manin-Mumford conjecture might need to be altered to include these exceptions. 

\subsection{Main Theorems}$ $

We introduce the projections involved in the setting of the universal Poincar\'e torsor:
\begin{align*}
&\Pi:\Pt\rightarrow \ag;\\
&\pi:\Pt\rightarrow \Ag\times_{\ag}\Agv;
\end{align*}
For the brevity of notations, we sometimes also denote $\Ag\times_{\ag}\Agv$ by $\Bg$.

We need another definition to state our main theorem, which seems strange but proves essential. This will be explained with more details in Proposition-Definition \ref{ribet_prop}.

\begin{definition}
    Let $S$ be a scheme, $A/S$ an abelian scheme with dual abelian scheme $\calA^\vee/S$. For a morphism of $S$-abelian schemes $\calA'\joinrel\xrightarrow{(\alpha,\beta)} \calA\times_S \calA^{\vee}$, if $\alpha^{\vee}\beta+\beta^{\vee}\alpha=0$, we say $\calA'$ (or $(\alpha,\beta)$) satisfies the \textbf{Ribet property}.
\end{definition}

In this article, we shall always use this definition when $\calA'$ is an $S$-abelian subscheme of $\calA\times_S \calA^\vee$, with $S$ being a quasi-projective variety over $\CC$ or $\QQ^{\mathrm{alg}}$. Here is the main classification result.

\begin{theorem}\label{thm_Pt}
For an irreducible subvariety $T\hookrightarrow\Pt$, let
$S=\Pi(T)\hookrightarrow \ag$.

Then, $T$ is a special subvariety of $\Pt$ if and only if the following holds:
\begin{enumerate}
    \item $S$ is a special subvariety of $\ag$ as pure Shimura varieties;
    \item One of the following holds:
    
    \begin{itemize}
        \item there is an abelian subscheme $\calA'$ of $(\Bg)_S$, and a torsion multi-section $t$, such that $T=\pi^{-1}(\calA'+t)$;
        \item there is an abelian subscheme $\calA'$ of $(\Bg)_S$ satisfying the \textbf{Ribet property} and $T$ is a torsion translate of the Ribet section (Definition \ref{ribet}) attached to $\calA'$. 
    \end{itemize}
\end{enumerate}
\end{theorem}

This theorem will be proved in Section \ref{section4}. An explanation of this classification will be given below Remark~\ref{RmkMSVSpecial}.

Through the natural projection $\Pt \rightarrow \Agv$, one can view $\Pt$ as the \textit{universal semi-abelian variety with toric part of rank $1$}. In general, we need to consider the universal semi-abelian variety with toric part of rank $n$, denoted by $\pgn$, for all $n\ge 1$. It can be constructed from $\Pt$ as follows. In the following, $X^{\times_{Y}n}$ means the $n$-th power fibered product \[\underbrace{X\times_{Y}X\times_{Y}\cdots \times_{Y}X}_{n \ \mathrm{times}}\] over a given structure morphism $X\to Y$. We also write $X^{[n]}$ when the base (and the structure morphism) is clear, especially for $\Ag\to \ag$ or $\Agv\to\ag$ (and thus defining $(\Ag)^{[n]}$, $(\Agv)^{[n]}$). Then $\pgn$ is defined by the following Cartesian diagram
\[\begin{tikzcd}
\pgn \arrow[rr] \arrow[d] \arrow[drr, phantom,"\ulcorner",very near start]                                               &  & (\Pt)^{\times_{S_g}n} \arrow[d]       \\
\Ag\times_{S_g}(\Agv)^{[n]} \arrow[rr, "{(\Delta_{\Ag},\mathrm{id})}"] &  & (\Ag\times_{S_g}\Agv)^{\times_{S_g}n}
\end{tikzcd}\]
Now $\pgn \rightarrow (\Agv)^{[n]}$ is the {\em universal semi-abelian variety with toric part of rank $n$} over the moduli space.

In other words, if we consider the natural projection $\pgn\rightarrow \Ag$, then relative to this base we have $$\pgn\cong\underbrace{\Pt\times_{\Ag}\Pt\times_{\Ag}\cdots \times_{\Ag}\Pt}_{n \ \mathrm{times}}.$$ %Then if $T$ is a special subvariety of $\pgn$, its projection to the $i$-th $\Pt$ is thus a special subvariety of $\Pt$, which has been described by Theorem \ref{thm_Pt}. Up to an isogeny of this semi-abelian scheme over $\Ag$, this completely determines the structure of the special subvariety $T$.

We have the following characterization of special subvarieties of $\pgn$, building upon Theorem~\ref{thm_Pt}.
\begin{theorem}\label{general}
    A special subvariety $T$ of $\pgn$ is an irreducible component of the inverse image of $T_1\times_{\Ag}T_2\times_{\Ag}\cdots\times_{\Ag}T_n$ under a matrix isogeny of $\pgn$ over $\Ag$, where the $T_i$'s are special subvarieties of $\Pt$ whose images in $\Ag$ are the same.
\end{theorem}
The term \textit{matrix isogeny} in the statement of the theorem refers to an isogeny of the form 
\[
\begin{tikzcd}
\underbrace{\Pt\times_{\Ag}\Pt\times_{\Ag}\cdots \times_{\Ag}\Pt}_{n \ \mathrm{times}} \arrow[rd] \arrow[rr, "{[h]}"] &     & \underbrace{\Pt\times_{\Ag}\Pt\times_{\Ag}\cdots \times_{\Ag}\Pt}_{n \ \mathrm{times}} \arrow[ld] \\
                                                     & \Ag &                                 
\end{tikzcd}\]for some $h\in\GL_n(\QQ)\cap \M_{n,n}(\ZZ)$.

\begin{remark}\label{RmkMSVSpecial}
Since any mixed Shimura variety whose pure part is a Hodge type Shimura variety is a family of semi-abelian varieties, the classification theorem above essentially also gives a classification of special subvarieties for all such mixed Shimura varieties.
\end{remark}






 We now briefly explain  what this classification in Theorem~\ref{thm_Pt} really says. In the universal family of abelian varieties over the Siegel moduli space, a special subvariety is essentially a translation of an algebraic subgroup scheme by a torsion (multi)-section (see \cite[Proposition 1.1]{gaoziyang}). However, in the case of the universal Poincar\'{e} torsor, these do not exhaust all special subvarieties. The main exceptional special subvarieties come from the Ribet sections. Note that a translation of an algebraic subgroup scheme by a torsion (multi)-section is a finite union of translations of algebraic subgroups on each fiber. The Ribet sections appear when this property holds fiberwise but not globally. More precisely, these are special subvarieties that are a finite union of translations of algebraic subgroups of $\Pt\rightarrow\Agv$ fiberwise (see Proposition \ref{fiber=subgp}) but not globally so.

Now we explain the precise meaning of the various translation processes appearing in the statements of Theorem \ref{thm_Pt}.
\begin{itemize}
    \item \textbf{The case of} $\Ag\times_{\ag}\Agv$.

    Let $S$ be a pure Shimura subvariety of $S_g$, whose $\CC$-points are of the form $S(\CC)=\Gamma_S\backslash X$. The restriction of the family $\Bg=\Ag\times_{\ag}\Agv$ to the base $S$ is a family of abelian varieties denoted by $(\Bg)_S$ in the theorem. Let $S'$ be a pure shimura variety whose $\CC$-points are of the form $S'(\CC)=\Gamma_{S'}\backslash X$ with $\Gamma_{S'}\subset\Gamma_{S}$. Then there is a natural finite morphism $S'\rightarrow S$ which on $\CC$-points is the natural projection $\Gamma_{S'}\backslash X\twoheadrightarrow \Gamma_{S}\backslash X$. A \textit{torsion multi-section} of the family $(\Bg)_S$ is a section $t$ of its base change $(\Bg)_{S'}=(\Bg)_S\times_S S'$ over $S’$, such that for all $x\in S'(\CC)$, the point $t(x)$ is a torsion point in the fiber over $x$. See the following diagram for an illustration.
\[\begin{tikzcd}
(\Bg)_{S'} \arrow[d] \arrow[r] \arrow[dr, phantom,"\ulcorner",very near start]     & (\Bg)_S \arrow[d] \\
S' \arrow[r] \arrow[u, "t", dotted, bend left] & S                    
\end{tikzcd}\]
    
    For a subvariety $Z$ of $(\Bg)_S$, take the base change $Z\times_S S'$, translate it by $t$, and then take an irreducible component of the image under the morphism $(\Bg)_{S'}\rightarrow (\Bg)_S$. The resulting subvariety of $(\Bg)_S$ is a \textit{torsion translation} of $Z$ and will be denoted by $Z+t$. This is the process of torsion translation on the level of subfamilies of $\Ag\times_{\ag}\Agv$.

    \item \textbf{The case of} $\Pt$.

    This case is slightly trickier since there are two group laws on $\Pt$ relative to the base $\Ag$ and $\Agv$ respectively. Again let $S$ be a pure Shimura subvariety of $\ag$ and $S'$ be a finite cover of it as in the case of $\Ag\times_{\ag}\Agv$. Then $(\Pt)_{S}$ is a family of semi-abelian varieties over $(\Agv)_S$. A \textit{torsion multi-section} of this family over it is a torsion section of $(\Pt)_{S'}$ over $(\Agv)_{S'}$ up to a self-isogeny. More precisely, it is a section denoted by $t$ in the following diagram
\[\begin{tikzcd}
f^{*}(\Pt)_{S'} \arrow[d] \arrow[r] \arrow[dr, phantom,"\ulcorner",very near start]                         & (\Pt)_{S'} \arrow[d] \arrow[r] \arrow[dr, phantom,"\ulcorner",very near start] & (\Pt)_S \arrow[d]  \\
(\Agv)_{S'} \arrow[r, "f"] \arrow[u, "t", dotted, bend left] & (\Agv)_{S'} \arrow[d] & (\Agv)_S \arrow[d] \\
                                                             & S' \arrow[r]                    & S                 
\end{tikzcd}\]
with $f$ being an isogeny, such that for all $x\in S'(\CC), y\in (\Agv)_x(\CC)$, the point $t(y)$ is a torsion point in the fiber over $y$.

 For a subvariety $Z$ of $(\Pt)_S$, take two base changes of it as in the above diagram, translate it by $t$, and then take an irreducible component of the image under the morphism 
 \[f^*(\Pt)_{S'}\rightarrow (\Pt)_{S'}\rightarrow (\Pt)_{S}.\] 
 The resulting subvariety of $(\Pt)_S$ is then a \textit{torsion translation} of it. Completely analogously, there is the notion of translation by a torsion multi-section of $(\Pt)_S$ over $(\Ag)_S$. A general torsion translation of a subvariety of $(\Pt)_S$ is then the result of a combination of these two translation processes. This is the process of torsion translation on the level of subfamilies of $\Pt$.
 
\end{itemize}


\subsection{Reformulation of the Relative Manin-Mumford Conjecture}$ $

Let the base $S$ be a regular, irreducible, quasi-projective variety over an algebraically closed field $k$ of characteristic $0$. Let us recall the Relative Manin-Mumford Conjecture for an abelian scheme over $S$, which is inspired by Shou-wu Zhang's ICM talk in 1998 \cite{Zhang1998} and formulated by Pink \cite{Pink05} and Zannier \cite{Zan12}. Gao and Habegger have recently proved this conjecture \cite{GaoHab_RMM}. There are a lot of works before this proof, notably when the subvariety in question is a curve by Masser and Zannier, see the references in \cite{GaoHab_RMM}. To make the statement closer to Pink's original formulation, we restate their theorem as follows.

\begin{definition}\label{DefnSpeSubAS}
    Let $\calA/S$ be an abelian scheme. A \textbf{special subvariety} of $\calA$ is the image of the translate of an abelian subscheme of $\calA\times_S S' \rightarrow S'$ by a torsion section, where $S'$ is a finite covering of a closed irreducible subvariety of $S$, under the natural projection $\calA\times_S S' \rightarrow \calA$.
\end{definition}


\begin{theorem}{\cite[Theorem 1.1]{GaoHab_RMM}}
    Let $\calA/S$ be an abelian scheme of relative dimension $g$, and $X$ be an irreducible closed subvariety of $\calA$. If $X(k) \cap \calA_{tor}$ is Zariski dense in $X$, then either $X$ is contained in a proper special subvariety of $\calA$, or $\dim X \geq g$.
\end{theorem}

%The notion of torsion translation is exactly the same as the discussions after Theorem \ref{general}, except that the base change is a general finite base change, rather than base change by a finite Shimura morphism. In fact, this is equivalent to translating by a torsion closed point in the geometric generic fiber $\calA_{\overline{\eta}}$, where $\overline{\eta}=\operatorname{Spec}(\overline{k(S)})$ is the spectrum of an algebraic closure of the function field of $S$. For a treatment in this spirit in the case of abelian schemes, we refer the reader to the discussion after the statement in \cite[Theorem 1.1]{GaoHab_RMM}.

Gao \cite[Proposition 1.1]{gaoziyang} classified all special subvarieties (in the sense of Pink) in $\Ag$, the universal abelian variety, with a natural projection $\pi:\Ag\to \ag$. He showed that an irreducible subvariety $Y$ of $\Ag$ is special if and only if the following holds:
    \begin{enumerate}
        \item $\pi(Y)$ is a special subvariety of $\ag$;
        \item $Y$ is, up to a finite covering of $\pi(Y)$, the translate of an abelian subscheme of $\pi^{-1}(\pi (Y))$ (over $\pi (Y)$) by a torsion section.
    \end{enumerate}
Definition~\ref{DefnSpeSubAS} can thus be translated in the language of moduli as follows. Let $\calA/S$ be an abelian scheme of relative dimension $g$. Up to a finite base change, we have the moduli map $S \rightarrow S_g$ and thus the moduli map $\calA \rightarrow \Ag$. Then an irreducible closed subvariety $X$ of $\calA$ is a special subvariety if and only if $X$ is an irreducible component of the pullback of a special subvariety of $\pi:\Ag\to \ag$ via the moduli map $\calA \rightarrow \Ag$.

%Then the moduli-theoretic description of Definition~\ref{DefnSpeSubAS}  is as follows:
%\begin{itemize}
%    \item Let $\calA/S$ be an abelian scheme. Let $X$ be an irreducible closed subvariety of $\calA$. Suppose $X$ is an irreducible component of the pullback of a special subvariety of $\pi:\Ag\to \ag$ via the classifying morphism $S\to \ag$ (up to a finite base change to gain a sufficiently high level structure). In this case, we say $X$ is \textbf{special}.
%\end{itemize}

The main results of our paper suggest that, for a family of semi-abelian varieties, the notion of special subvarieties should be defined similarly, and is suitable for the Relative Manin-Mumford Conjecture. Here is our definition. Let $\calG/S$ be a semi-abelian scheme of relative dimension $g+n$, with toric part of rank $n$. Up to a finite \'{e}tale base change (to gain a sufficiently high level structure), we get the moduli map $S\to (\Agv)^{[n]}$ and hence the moduli map $\calG\rightarrow \pgn$. %Accept for the moment that, given a Shimura family $\calG\to S$ of semi-abelian varieties (meaning that $\calG$ and $S$ are Shimura varieties, and the morphism is a Shimura morphism.), then every semi-abelian subvariety over a Hodge generic point/ $s$ of $S(\CC)$ can be extend to a semi-abelian scheme over whole of $S$.\footnote{The readers familiar with Hodge structures can prove this easily. The only thing to do is to consider the orbit of the point under its Mumford-Tate group, and use the period map.} 

\begin{definition}
    Let $Z$ be an irreducible closed subvariety of $\calG$. We say that $Z$ is \textbf{special}, if $Z$ is an irreducible component of the pullback of a special subvariety of $\pgn$ via the moduli map $\calG \rightarrow \pgn$.% In this case, we say $Z$ is \textbf{special}.
\end{definition}

By Theorem \ref{general}, such special subvarieties are essentially combinations of closed subgroup schemes and Ribet sections, which seems to be the first step towards the following reformulated conjecture:

\begin{conjecture}[Relative Manin-Mumford Conjecture for semi-abelian schemes]$ $ \label{intro_RMM}
    Let $\calG/S$ be a semi-abelian scheme of relative dimension $g+n$, with toric rank $n$. Let $Z$ be an irreducible closed subvariety of $\calG$. If $Z(k) \cap \calG_{tor}$ is Zariski dense in $Z$, then either $Z$ is contained in a proper special subvariety of $\calG$, or $\dim X \geq g+n$.
\end{conjecture}

%A more detailed exposition of the meaning of the inequality “$\dim X\geq g+n$” in the context of unlikely intersections is given 
In Section \ref{SubsectionReformulationRMM}, we will prove that Conjecture~\ref{intro_RMM} is a corollary of the Zilber-Pink Conjecture for mixed Shimura varieties.


\subsection{Zilber-Pink for mixed Shimura varieties implies Zilber-Pink for semi-abelian varieties}
In Section~\ref{SubsectionZPimpliesZPSA}, we will give a proof that the Zilber-Pink Conjecture for mixed Shimura varieties implies the Zilber-Pink Conjecture for semi-abelian varieties, \textit{i.e.} \cite[Conjecture~1.3]{Pink05} implies \cite[Conjecture~5.1]{Pink05}. This was done by Edixhoven \cite{EdixhovenPrivate}, as announced in \cite[Remark~1.3]{BE20}. Both our treatments are based on the following structural result: for any special subvariety $T$ of the universal semi-abelian variety $\pgn \rightarrow (\Agv)^{[n]}$, the fiber $T_s$ is a finite union of translates of subgroups of $(\pgn)_s$ for any $s \in (\Agv)^{[n]}(\mathbb{C})$. We prove this as Proposition~\ref{fiber=subgp} and Edixhoven proves a finer veresion as \cite[Proposition~3.9]{EdixhovenPrivate}. However, our proofs to this structural result are different: we did this as an application of our  geometric characterization of special subvarieties of $\pgn$ (Theorem~\ref{general}), while Edixhoven did this by a direct group-theoretic computation based on and generalizing \cite{BE20} and without going into the full geometric structure of special subvarieties of $\pgn$.


\subsection{Structure of the Article}$ $

The structure of this article is as follows. In Section \ref{section2} we introduce the concepts of Ribet property and Ribet sections, generalizing the ones in \cite{BE20}. They are the central object of the article. Then we check Ribet sections are special subvarieties of $\Pt$ in Section \ref{section3}. This section mainly focuses on the mixed Hodge structures and can be skipped on the first reading, and the computation largely follows the treatment of \cite{BE20}. 
 Section \ref{section4} and \ref{section5} are another core of this article, where we prove our classification results Theorem~\ref{thm_Pt} and Theorem~\ref{general}. %To help readers understand the geometric meaning, the full understanding of the main results and the moving step (see Subsection \ref{moving}). 

Finally in Section \ref{section6}, we give a more plausible reformulation of the Relative Manin-Mumford Conjecture and collect the implications of the conjectures involved. We will also prove that Zilber-Pink for mixed Shimura varieties implies Zilber-Pink for semi-abelian varieties. Both are based on our main result Theorem~\ref{general}. %On the first reading, the statements and the implication relations of conjectures seem to be of most interest. We also hope this section will serve as an introduction to this circle of ideas. $ $

\subsection*{Acknowledgements}
This work is one of the projects in the summer school “Algebra and Number Theory” held by Peking University and the Academy of Mathematics and Systems Science, CAS. We want to express our gratitude to Professor Shou-wu Zhang and Professor Liang Xiao. Without their efforts, the summer school could not held with such a success. We are also grateful to our mentor, Professor Ziyang Gao, for suggesting this question and continuous guidance. 
We thank Professor Liang Xiao and Professor Xinyi Yuan for their interest in the results. We also thank our group members Houjun Chen and Ziheng Zhang for their participation and discussion.
CH would like to thank Peking University for its great hospitality.
KG would like to thank Wenbin Luo for his guidance on \LaTeX writing.
\section{Construction of Ribet Sections}\label{section2}

In this section, we will generalize the construction of Ribet sections in \cite{BE20}. Let $S$ be a scheme, and let $\calA$ be an abelian scheme over $S$. The dual abelian scheme is denoted by $\calA^{\vee}$ (for a reference of the notion of the dual of an abelian scheme and several constructions for abelian schemes, see \cite{faltings-chai}). Let $\pl_\calA$ be the Poincar\'{e} line bundle over $\calA\times_{S}\calA^{\vee}$ with trivializations along $0_\calA\times_S \calA^{\vee}$ and $\calA\times_S 0_{\calA^{\vee}}$, compatible over $0_\calA\times_S 0_{\calA^\vee}$. When restricted to certain abelian subschemes of $\calA\times_{S}\calA^{\vee}$, the Poincar\'{e} line bundle would become algebraic equivalent to the trivial one. In fact, it would admit a trivialization after a finite cover in such case, through which the Ribet (multi)-section is defined. To demonstrate this, we have the following equivalence at the first step.

%This notion may differ from that found in the literature, but it is convenient for our purposes. If $\calA$ admits a principal polarization, the isotropic subschemes will have a concrete description, which supports the usage of the terminology “isotropic”. (origin explanation of ribet_prop)

\begin{prop-def}\label{ribet_prop}
Let $S$ be a connected, reduced noetherian scheme, and $\calA/S$ an abelian scheme over $S$. For an $S$-abelian subscheme $\calA'\stackrel{(\alpha,\beta)}{\longrightarrow} \calA\times_S \calA^{\vee}$, the followings are equivalent:

(1) $(\alpha,\beta)^*\pl_\calA\in \pic^0_{\calA'/S}(S)$;

(2) $\alpha^{\vee}\beta+\beta^{\vee}\alpha=0$;

(3) $((\alpha,\beta)^*\pl_\calA)^{\otimes 2}$ is pulled back from $S$ (by abuse of notation, we also say such a line bundle is trivial).

If one of the conditions above holds for $(\alpha,\beta):\calA'\hookrightarrow \calA\times_S\calA$ over $S$, we say $\calA'$ (or $(\alpha,\beta)$) satisfies the \textbf{Ribet property}.
\end{prop-def}

By abuse of notation, we also denote by $\pl_\calA|_{\calA'} := (\alpha,\beta)^*\pl_\calA$.

\begin{proof}
    Notice that (3) implies (1). For any line bundle $\calM$ on $\calA'/S$, it lies in $\pic^0_{\calA'/S}(S)$ if and only if $\phi_\calM=0$. Here the notation $\phi_\calM$ denotes the $\phi$-construction. Recall that it is a homomorphism from $\calA$ to $\calA^{\vee}$, such that for any $S$-scheme $T$, the homomorphism is given on $T$-points by
    \begin{align*}
        \phi_{\mathcal{M}}: \calA(T)&\rightarrow \calA^{\vee}(T)\\
        x&\mapsto t_x^{*}(\mathcal{M}_{\calA_T})\otimes \mathcal{M}_{\calA_T}^{-1}.
    \end{align*}
    Moreover, the homomorphism $\phi_\lb$ is an isogeny if $\calM$ is ample or anti-ample. For a reference of these facts, see \cite{Mum12}.

    
    Thus, if $\calM^{\otimes2}$ is trivial, then $0=\phi_{\calM^{\otimes2}}=2\phi_\calM$. Cancelling the isogeny, we get $\phi_\calM=0$, which implies $\calM\in \pic^0_{\calA'/S}(S)$.

\begin{enumerate}
    \item 
    First, assume that $S=\spec(k)$ is the spectrum of an algebraically closed field. Choose an ample line bundle $\mathcal{L}$ on $\calA$ and choose an isogeny $\psi: \calA^{\vee}\rightarrow \calA$ with $\phi_{\mathcal{L}}\circ\psi=[N]_{\calA'}$ for a positive integer $N$.

    Let $\gamma:=\psi\circ \beta$, and $m:\calA\times_S \calA \to \calA$ be the multiplication. Note that the composition $\calA'\joinrel\xrightarrow{(\alpha, \gamma)} \calA\times_S \calA\joinrel\xrightarrow{(\id,\phi_{\mathcal{L}})} \calA\times_S \calA^\vee$ is  $(\alpha, N\beta)$. It is a standard fact that there is a canonical isomorphism $(\id,\phi_{\mathcal{L}})^*\pl_\calA=\Lambda(\lb)=m^*\lb\otimes \pr_1^*\lb^{\otimes-1}\otimes\pr_2^*\lb^{\otimes-1}$. Thus, we obtain an explicit description: $$(\alpha, N\beta)^*\pl_\calA=(\alpha+\gamma)^*\lb\otimes\alpha^{*}\lb^{\otimes-1}\otimes\gamma^{*}\lb^{\otimes-1}.$$ For the simplicity of notations, let $$\calM_{(\alpha,\gamma)}(\lb):=(\alpha+\gamma)^*\lb\otimes\alpha^{*}\lb^{\otimes-1}\otimes\gamma^{*}\lb^{\otimes-1}.$$
    There is a criterion for such a line bundle to be in $\pic^0(\calA')$.

    \begin{lemma} \label{M_in_pic0}
        Let $\gamma:=\psi\circ \beta$ as before. Then, $\alpha^{\vee}\beta+\beta^{\vee}\alpha=0$ is equivalent to $\calM_{(\alpha,\gamma)}(\lb)\in \pic^0(\calA')$.
    \end{lemma}

    \begin{proof}
        We compute \begin{align*}
            \phi_{\calM_{(\alpha,\gamma)}(\lb)}&=\phi_{(\alpha+\gamma)^*\lb}-\phi_{\alpha^*\lb}-\phi_{\gamma^*\lb}\\
            &=(\alpha+\gamma)^\vee\circ\phi_\lb\circ (\alpha+\gamma)-\alpha^\vee\circ\phi_\lb\circ\alpha-\gamma^\vee\circ\phi_\lb\circ\gamma
            \\&=\alpha^\vee\circ\phi_\lb\circ\gamma+\gamma^\vee\circ\phi_\lb\circ\alpha
            \\&=\alpha^\vee\circ[N]_{\calA'}\circ\beta+\beta^\vee\circ[N]_{\calA}\circ\alpha
            \\
            &=(\alpha^\vee\beta+\beta^\vee\alpha)\circ [N]_{\calA'}.
        \end{align*}
        Thus the equivalence follows.
    \end{proof}

    To see the relation between $(\alpha,N\beta)^*\pl_\calA$ and $(\alpha,\beta)^*\pl_\calA$, we also need the following result:
    \begin{lemma}\cite[Exercise (7.4)]{BGB_AV}
        For $m,n\in \ZZ$, the pullback of Poincar\'e line bundle on $\calA\times_S \calA^\vee$ via $(m,n):\calA\times_S \calA^\vee \to \calA\times_S \calA^\vee$ is given by $(m,n)^*\calP_\calA\cong \calP_\calA^{\otimes mn}$.
    \end{lemma}
    
    \begin{proof}
        It suffices to show $(m,1)^*\calP_\calA\cong \calP_\calA^{\otimes m}$, because of the identity $(m,n)=(m,1)\circ(1,n)$ and the self-duality of abelian varieties. 
        
        View $\calA\times_S \calA^\vee\to \calA^\vee$ as a family of abelian varieties over $\calA^\vee$, with the universal family $\calP_\calA$ of algebraically trivial line bundles. Thus $[\calP_\calA] \in \pic^0_{\calA\times_S \calA^\vee/\calA^\vee}(\calA^\vee)$. By the linearlity of pullbacks on $\pic^0$, we have $[(m,1)^{*}\calP_\calA]=[\calP_\calA^{\otimes m}]$ as an equivalence class in $ \pic^0_{\calA\times_S \calA^\vee/\calA^\vee}(\calA^\vee)$. Use the rigidification of $\calP_\calA$ along $0_\calA\times_S \calA^\vee$, one obtains $(m,1)\calP_\calA\cong\calP_\calA^{\otimes m}$ as line bundles on $\calA\times_S \calA^\vee$.
    \end{proof}

    Note the following diagram, where the equality of line bundles over $\calA\times_S \calA^\vee$ is the content of the above lemma:
    
    \begin{center}
\begin{tikzcd}
                                                                                                                                                                                                           &                                                & \Lambda(\lb)^{\otimes N} \arrow[d, no head] & {(N,1)^*(\calP_\calA)=\calP^{\otimes N}_\calA} \arrow[d, no head] & \calP_\calA \arrow[d, no head] \\
\calA' \arrow[r, "{(\alpha,\beta)}"'] \arrow[rr, "{(\alpha,\phi_\lb\circ \beta)=(\alpha,\gamma)}", bend left] \arrow[rrrr, "{N\circ(\alpha,\beta)}"', bend right] \arrow[rrr, "{(\alpha,N\beta)}", bend right] & \calA\times_S \calA^\vee \arrow[r, "{(\id, \phi_\lb)}"'] & \calA\times_S \calA \arrow[r, "{(\id,\psi)}"']        & \calA\times_S \calA^\vee \arrow[r, "{(N,1)}"]                       & \calA\times_S \calA^\vee            
\end{tikzcd}
    \end{center}

    The line bundles are all pulled back from $\calP_\calA$ on the rightmost of the diagram. Compare the resulting line bundles on $\calA'$, one obtains the equality: 
    $$(\alpha,\beta)^*(\pl_\calA^{\otimes N^2})=(\alpha, N\beta)^*(\pl_\calA^{\otimes N})=\calM_{(\alpha,\gamma)}(\lb)^{\otimes N}.$$
    Thus, (2) implies (1). Indeed, (2) implies $\calM_{(\alpha,\gamma)}(\lb)\in \pic^0(\calA')$. Further, one uses above equality to conclude $(\alpha,\beta)^*\pl_\calA\in \pic^0(\calA')$, since a non-zero power of it lies in $\pic^0(\calA')$. (Here's a quick proof of this fact. For any non-zero integer $M$, $\phi_{\lb^{\otimes M}}=M\phi_{\lb}=0$ is equivalent to $\phi_{\lb}=0$.)
    
    Now we show (2) implies (3). Here comes the trick. From the above observations, if $\alpha^\vee\beta+\beta^\vee\alpha=0$, then
    $$(\pl_\calA|_{\calA'})^{\otimes 4}\cong [2]_{\calA'}^*\pl_\calA|_{\calA'}\cong(\pl_\calA|_{\calA'})^{\otimes 2}.$$
    Cancelling out, we obtain that $(\pl_\calA|_{\calA'})^{\otimes2}$ is trivial.

    \item Now we prove the statement for a general base $S$ as in the statement.
    
    We show (2) implies (3). Assume (2) now. Since the formation of $(\pl_\calA|_{\calA'})^{\otimes 2}$ is compatible with base changes of $S$, one concludes $(\pl_\calA|_{\calA'_{\bar s}})^{\otimes 2}$ is trivial for any geometric fiber $\calA'_{\bar s}$ of $\calA'\to S$ for a geometric point $\bar s\rightarrow S$ from previous results. By the Seesaw theorem, there is a maximal closed subscheme $K\hookrightarrow S$, such that $(\pl_\calA|_{\calA'})^{\otimes 2}$ is trivial over $K$ (i.e. its restriction on $\calA'_K$ is pulled back from $K$, and $K$ is universal among this property). The maximality of $K$ then implies every geometric point $\bar s\to S$ factors through $K$, and thus $K=S$ since $S$ is reduced.

    The implication (1) $\Rightarrow$ (2) proceeds similarly. For every geometric point $\bar s\to S$, ${{\pl_{\calA}} |_{\calA'_{\bar s}}} \in \pic^0(\calA'_{\bar s})$ implies $(\alpha^\vee\beta+\beta^\vee\alpha)_{\bar s}=\alpha_{\bar s}^\vee\beta_{\bar s}+\beta_{\bar s}^\vee\alpha_{\bar s}=0$ again by \ref{M_in_pic0}. Thus the homomorphism $\alpha^\vee\beta+\beta^\vee\alpha$ is itself zero as it is zero on every (geometric) fiber of $\calA'\rightarrow S$.
\end{enumerate}
\end{proof}

Given this result, we can imitate the constructions in \cite{BE20}. Consider the functor $\underline{\mathrm{Isom}}(\mathcal{O}, \pl_\calA)$:
\begin{align*}
     \mathsf{Sch}_{\calA\times_S \calA^{\vee}}&\rightarrow \mathsf{Set}\\
    [T\stackrel{f}{\rightarrow}\calA\times_S \calA^{\vee}]&\mapsto\mathrm{Isom}_{\mathrm{Mod}_{\mathcal{O}_T}}(\mathcal{O}_T,f^*\pl_\calA).
\end{align*}
This functor is represented by a scheme $\pt_\calA$ over $\calA\times_S \calA^{\vee}$ (the geometric line bundle of $\pl_\calA$ with the zero section removed), which is a $\gm$-torsor over $\calA\times_S \calA^{\vee}$. Upon the two projections from $\calA\times_S \calA^{\vee}$ to $\calA$ and $\calA^{\vee}$, the Poincar\'{e} torsor $\pt_\calA$ is realized as a family of semi-abelian schemes over $\calA$ and $\calA^\vee$, by which it is also called the Poincar\'{e} bi-extension in the literature.

Now we introduce the construction of Ribet sections. Let $\calA'$ be an abelian subscheme of $\calA\times_S \calA^\vee$ satisfying the Ribet property. Consider the following diagram:
\[
\begin{tikzcd}
{[2]_{\calA'}^*(\pt_\calA|_{\calA'})} \arrow[d] \arrow[r, "\pi"]\arrow[dr, phantom,"\ulcorner",very near start]            & \pt_\calA|_{\calA'} \arrow[d] \arrow[r]\arrow[dr, phantom,"\ulcorner",very near start] & \pt_\calA \arrow[d]    \\
\calA' \arrow[r, "{[2]_{\calA'}}"] \arrow[u, "\tilde{r}_{\calA'}", dotted, bend left] & \calA' \arrow[r, hook]              & \calA\times_S \calA^{\vee}
\end{tikzcd}
\]
By Proposition \ref{ribet_prop}, $[2]_{\calA'}^*(\pl_\calA|_{\calA'})\cong (\pl_\calA|_{\calA'})^{\otimes 2}$ is trivial. Moreover, since the Poincar\'{e} line bundle is bi-rigidified, this trivialization is unique. Hence by the definition of $\pt_\calA$, this determines a section $\tilde{r}_{\calA'}$ of $[2]_{\calA'}^*(\pt_\calA|_{\calA'})\rightarrow \calA'$. 
\begin{definition}\label{ribet}
The image of the composite $r_{\calA'}:=\pi\circ \tilde{r}_{\calA'}$ is called the \textbf{Ribet (multi-)section} over $\calA'$.
\end{definition}

\begin{remark}
    In \cite{BE20}, the authors considered a special kind of abelian subschemes with the Ribet property, and their notion of the Ribet section is a special case of ours. Let $f:\calA\rightarrow \calA^{\vee}$ be a homomorphism of abelian schemes and $\alpha=f-f^{\vee}$. Up to a switch of factors, their construction essentially deals with the subscheme $\calA'=\Gamma(\alpha)$, the graph of $\alpha$. In this case, the authors verified a stronger version of the Ribet property by showing that $\pl_\calA|_{\calA'}$ is already trivial. Hence there's no need to do base change by $[2]_{\calA'}$ like above, and the Ribet section is genuinely a section, rather than merely a multi-section. 
\end{remark}
Reflecting on the definition of the Ribet property, one may notice the analogy of an abelian subvariety with the Ribet property and an isotropic subspace in a symplectic space, with the Poincar\'{e} bundle playing the role of a symplectic form. Recall that in a non-degenerate symplectic space of dimension $2g$, an isotropic subspace has dimension $\leq g$, with the maximum attained by Lagrangian subspaces. We have the following proposition stating that this dimension bound actually holds for an abelian subvariety with the Ribet property. The proof is easy, but the analogy it reveals is important.

\begin{proposition}
Let $A$ be an abelian variety over a field $k$ of dimension $g\geq 1$, and let $A'\subset A\times A^\vee$ be an abelian subvariety with the Ribet property. Then the dimension bound
\[\dim A'\leq g\]
holds.
\end{proposition}

\begin{proof}
Prove by contradiction. Suppose $A$ is an abelian variety over a field $k$ of dimension $g\geq 1$, and $A'\subset A\times A^\vee$ is an abelian subvariety with the Ribet property of dimension $>g$. Choose an ample line bundle $\lb$ on $A$ and let $\phi_\lb:A\rightarrow A^\vee$ be the polarization constructed using the $\phi$-construction. Let $\Gamma(\phi_\lb)\subset A\times A^{\vee}$ be the graph of $\phi_\lb$. Finally, let $A''=(\Gamma(\phi_\lb)\cap A')^{\circ}$ be the identity component of $\Gamma(\phi_\lb)\cap A'$, which is a non-trivial abelian subvariety of $A\times A^\vee$.

Consider the restriction of the Poincar\'{e} line bundle $\calP_A$ to $A''$. On one hand, $\calP_A|_{A''}$ is in $\operatorname{Pic}^0(A'')$ since $A'\subset A\times A^\vee$ is an abelian subvariety with the Ribet property. On the other hand, $\calP_A|_{A''}$ is ample since $\calP_A|_{\Gamma(\phi_{\lb})}$ is ample. These two observations lead to a contradiction.
\end{proof}
The following immediate consequence will be used in Section \ref{section6}.
\begin{corollary}\label{not_entire}
Let $A$ be an abelian variety over a field $k$ of dimension $g\geq 1$, and let $A'\subset A\times A^\vee$ be an abelian subvariety with the Ribet property. Then $A\times \{s\} \subset A'$ for a closed point $s$ in $A^\vee$ if and only if $A'=A\times 0_{A^{\vee}}$ and $s=0_{A^\vee}$.
\end{corollary}

In \cite{BE20}, the authors verified that if the base scheme $S$ is a certain Shimura subvariety of the Siegel modular variety, the special kind of Ribet sections they constructed are special subvarieties of $\pt_\calA$ as a mixed Shimura variety. In the next section, we will be in the same context to verify that the general Ribet sections we constructed are special subvarieties.

\section{Ribet Sections as Special Subvarieties}\label{section3}

In this section, we review some constructions in mixed Hodge theory to show that the Ribet sections we constructed in Section \ref{section2} are special subvarieties of the universal Poincar\'{e} torsor. Our treatment largely follows \cite[Section 4-5]{BE20}.
\subsection{Review of Mixed Hodge Structures}\label{MHS}
\subsubsection{Pure Hodge Structures}$ $

For $n\in\mathbb{Z}$, a \textit{$\ZZ$-Hodge structure of weight $n$} is a finitely generated $\ZZ$-module with a decomposition
\[M_{\CC}=\bigoplus_{p+q=n}M^{p,q}\]
such that for all $p,q\in\ZZ$ with $p+q=n$, $\overline{M^{p,q}}=M^{q,p}$. There are standard notions like morphisms between pure Hodge structures, $\Hom$'s, and tensor products of Hodge structures, which we do not repeat here. The \textit{Tate structure} $\ZZ(n)$ is a Hodge structure of weight $-2n$, whose underlying module is the $\ZZ$-submodule $(2\pi i)^n\ZZ$ of $\CC$, and whose Hodge decomposition is set to be $\ZZ(n)_\CC=\ZZ(n)^{-n,-n}$, where $\ZZ(n)^{-n,-n}\cong \CC$ is the only non-zero component.

Consider the $\CC$-linear operator $C$ on $M_{\CC}$ which is multiplication by $i^{-p}\overline{i}^{-q}$ on $M^{p,q}$. It is easy to check that it descends to an $\RR$-linear operator on $M_{\RR}$ which we still denote by $C$. A \textit{polarization} on $M$ is a morphism of Hodge structures $\Psi: M\otimes M\rightarrow \ZZ(-n)$ such that the $\RR$-bilinear form
\[M_\RR\times M_\RR\rightarrow\RR,\ (x,y)\mapsto(2\pi i)^n\Psi(x\otimes Cy)\]
is symmetric and positive definite.
\begin{example}[Principally polarized abelian varieties]\label{ppav}$\ $

Let $g\geq 1$ be an integer. Consider a free $\ZZ$-module $M$ of rank $2g$ with a Hodge decomposition $M_{\CC}=M^{-1,0}\oplus M^{0,-1}$ and a polarization $\Psi:M\otimes M\rightarrow \ZZ(1)$ which induces an isomorphism of Hodge structures $M\rightarrow M^{\vee}(1)$. From this polarized Hodge structure, one can construct a principally polarized abelian variety of dimension $2g$, whose $\CC$-points are given by $M\backslash M_{\CC}/M^{0,-1}$. It was a theorem of Riemann that this construction is an equivalence of categories.

One can choose a basis of $M$ such that $M$ together with $\Psi$ is isomorphic to $\ZZ^{2g}$ with
\[\Psi:\ZZ^{2g}\otimes \ZZ^{2g}\rightarrow \ZZ(1),\ x\otimes y\mapsto 2\pi i x^\top\begin{pmatrix}\begin{smallmatrix} 0&-1_g\\1_g&0
\end{smallmatrix}\end{pmatrix}y.\]
The set $D_{\Psi}$ of Hodge structures of type $\{(-1,0),(0,-1)\}$ on $\ZZ^{2g}$ for which $\Psi$ is a polarization is in bijection with the Siegel upper half space $\HH_{g}:=\{\tau\in\M_{g}(\CC):\tau^{\top}=\tau,\ \mathrm{Im}(\tau)>0\}$ via
\[\tau\mapsto M^{0,-1}_{\tau}:=\begin{pmatrix} \tau\\1_g
\end{pmatrix}\CC^g.\]
\end{example}

\begin{remark}\label{poltype}
For a general $\Psi: M\otimes M\rightarrow \ZZ(1)$ which is not necessarily a principal polarization, one can choose a basis of $M$ such that $M$ together with $\Psi$ is isomorphic to $\ZZ^{2g}$ with
\[\Psi:\ZZ^{2g}\otimes \ZZ^{2g}\rightarrow \ZZ(1),\ x\otimes y\mapsto 2\pi i x^\top\begin{pmatrix}\begin{smallmatrix} 0&-\delta\\\delta&0
\end{smallmatrix}\end{pmatrix}y.\] 
Here $\delta=\mathrm{diag}\{e_1,...,e_g\}$ for positive integers $e_1,...,e_g$ with $e_1|e_2|...|e_g$. In this case, to make the above example work, one simply changes the Siegel upper half space $\HH_g$ by 
\[\delta^{-1}\HH_g:=\{\tau\in\M_{g}(\CC):(\delta\tau)^{\top}=\delta\tau,\ \mathrm{Im}(\delta\tau)>0\}.\] All contents in the entire article still hold if one replaces all $\HH_g$ by $\delta^{-1}\HH_g$. For moduli interpretations, this just replaces the moduli problem of principally polarized abelian varieties with the moduli problem of polarized abelian varieties of polarization type $\delta$. For the simplicity of notations, we simply deal with the case of principally polarized abelian varieties in all group-theoretical computations in this article. But note that the statement of the Ribet property (Proposition-Definition \ref{ribet_prop}) does not need the existence of a principal polarization.
\end{remark}

\subsubsection{Mixed Hodge Structures}$ $

A \textit{$\ZZ$-mixed Hodge structure} is a finitely generated $\ZZ$-module $M$ with the following structures. 
\begin{itemize}
\item There is an increasing filtration $\{W_n M\}_{n\in\ZZ}$ on $M$ called the \textit{weight filtration} such that $W_n M=M_{\text{tor}}$ for $n\ll 0$ and $W_n M=M$ for $n\gg 0$ with all $M/W_nM$ torsion free. 
\item There is a decreasing filtration $\{F^pM_\CC\}_{p\in\ZZ}$ on $M_\CC$ called the \textit{Hodge filtration} such that $F^p M=0$ for $p\gg 0$ and $F^pM=M$ for $p\ll 0$. The filtration induced by $F$ on $(\Gr^W_n M)_{\CC}$ should make a Hodge structure of weight $n$:
\[(\Gr^W_n M)_{\CC}=\bigoplus_{p+q=n}(\Gr^W_n M)_{\CC}^{p,q}\]
with 
\[(\Gr^W_n M)_{\CC}^{p,q}=F^p(\Gr^W_n M)_{\CC}\cap\overline{F^q(\Gr^W_n M)_{\CC}}.\]
\end{itemize}
\begin{example}[A type of 1-motives related to the Poincar\'{e} torsor]$\ $

In this example we will illustrate the mixed Hodge structures considered in \cite[Subsection 4.4]{BE20} using the language of 1-motives in \cite{hodge3}. We find this language helpful to understand the moduli interpretations of the universal Poincar\'{e} torsor as a mixed Shimura variety. This language will also be instructive when we construct the mixed Shimura data of Ribet sections. 

First, we recall the definition of 1-motives.
\begin{definition}{\cite[D\'{e}finition 10.1.2]{hodge3}.}$\ $

    A \textbf{1-motive} over an algebraically closed field $k$ is the data consist of 
    \begin{itemize}
    \item A free abelian group $X$ and a semi-abelian variety $G$ over $k$;
    \item A homomorphism of groups $X\stackrel{u}{\rightarrow}G(k)$.
    \end{itemize}
    A 1-motive given by these data is denoted by $[X\stackrel{u}{\rightarrow}G]$.
\end{definition}

Let $g\geq 1$ be an integer and $M=\ZZ(1)\oplus \ZZ^{2g}\oplus \ZZ$ with standard basis $2\pi i e_0,e_1,...,e_{2g+1}$. Impose a weight filtration on $M$ as follows:
\begin{align*}
W_{-3}M&=\{0\},\\
W_{-2}M&=\ZZ2\pi i e_0,\\
W_{-1}M&=\ZZ2\pi i e_0\oplus \ZZ e_1\oplus\cdots\oplus \ZZ e_{2g},\\
W_{0}M&=M.
\end{align*}
Let $D$ be the set of decreasing filtrations $F^{\bullet}$ on $M_{\CC}$ such that $(M,W_{\bullet},F^{\bullet})$ is a $\ZZ$-mixed Hodge structure of type $\{(-1,-1),(-1,0),(0,-1),(0,0)\}$ and such that 
\[\Psi:\ZZ^{2g}\otimes \ZZ^{2g}\rightarrow \ZZ(1),\ x\otimes y\mapsto 2\pi i x^\top\begin{pmatrix}\begin{smallmatrix} 0&-1_g\\1_g&0
\end{smallmatrix}\end{pmatrix}y\]is a polarization of pure Hodge structure on $\Gr^{W}_{-1}M$ under the isomorphism $\Gr^{W}_{-1}M\cong \ZZ^{2g}$ given by the chosen basis. It is easy to see that a Hodge filtration in $D$ is determined by $F^0$, that is, it is of the form
\[M=F^{-1}M_\CC\supset F^0M_{\CC}\supset F^{1}M_\CC=0.\]
\end{example}
\begin{proposition}\label{D}{\cite[Proposition 4.5]{BE20}}$\ $

There is a bijection $\HH_g\times \M_{1,g}(\CC)\times \M_{g,1}(\CC)\times \CC \mapsto D$,
\[(\tau,u,v,w)\mapsto F^0 M_{\CC}=\begin{pmatrix}
    u&w\\
    \tau& v\\
    1_g&0\\
    0&1
\end{pmatrix}\CC^{g+1}.\]
\end{proposition}

Now consider a mixed Hodge structure $(M,W_{\bullet},F^{\bullet})$ on $M$ given by a filtration $F^{\bullet}$ in $D$. One can define a 1-motive over $\CC$ by 
\[[W_{-1}M\backslash W_0M\rightarrow W_{-1}M\backslash M_{\CC}/F^0M_{\CC}].\] 
Here $W_{-1}M\backslash W_0M\cong\ZZ$ and $W_{-1}M\backslash M_{\CC}/F^0M_{\CC}$ is a semi-abelian variety whose abelian variety quotient is the principally polarized abelian variety \[(\Gr^W_{-1}M)\backslash (\Gr^W_{-1}M)_{\CC}/F^0(\Gr^{W}_{-1}M)_{\CC}.\]
In other words, this amounts to taking a semi-abelian variety over $\CC$ which is an extension of a principally polarized abelian variety by $\gm$, and taking a $\CC$-point on it.

Recall that for $A$ an abelian variety over $\CC$, there is a canonical isomorphism $A^{\vee}(\CC)\cong \operatorname{Ext}_{\mathsf{CommAlgGrp}_{/\CC}}(A,\gm)$. Therefore the moduli space of the semi-abelian variety part of such 1-motives should be the \textit{dual} of the universal abelian variety over the moduli space of principally polarized abelian varieties, and the moduli space of such 1-motives should be the universal semi-abelian variety over it. In other words, this is the universal Poincar\'{e} torsor of the universal abelian variety over the moduli space of principally polarized abelian varieties.

Now we are on the stage to illustrate the main computation in \cite[Section 4]{BE20}. Let $P$ be the subgroup scheme of $\GL(M)\times \GL(\ZZ(1))$ that fixes 
\begin{itemize}
    \item $W_{\bullet}$, 
    \item $\ZZ(1)\rightarrow W_{-2}M$, $2\pi i a\mapsto 2\pi i a e_0$, 
    \item $\ZZ\rightarrow \Gr_0^{W}M$, $a\mapsto ae_{2g+1}$,
    \item $\Psi:\Gr^{W}_{-1}M\otimes \Gr^{W}_{-1}M\rightarrow \ZZ(1)$. 
\end{itemize}
\begin{lemma}
    The $R$-valued points of $P$ for a ring $R$ are
\begin{align*}
    P(R)=\left\{\begin{pmatrix}
    \mu(g)&x&z\\
    0&g&y\\
    0&0&1
\end{pmatrix}:
    (g,\mu(g))\in\GSp(\Psi)(R),
x\in\M_{1,2g}(R),y\in\M_{2g,1}(R),z\in R 
 \right\}.
\end{align*}
\end{lemma}
The $R$-valued points of its unipotent radical are
\begin{align*}
    P^u(R)=\left\{\begin{pmatrix}
    1&x&z\\
    0&1&y\\
    0&0&1
\end{pmatrix}:
x\in\M_{1,2g}(R),y\in\M_{2g,1}(R),z\in R 
 \right\}.
\end{align*}
And $P$ has a normal subgroup $U$, whose $R$-valued points are
\begin{align*}
    U(R)=\left\{\begin{pmatrix}
    1&0&z\\
    0&1&0\\
    0&0&1
\end{pmatrix}:z\in R 
 \right\}.
\end{align*}

We refer to \cite[Section 4]{BE20} for some explicit formulas of the actions of these groups on $D$. Quotienting of $D$ by the action of $P(\ZZ)$ should be viewed as eliminating the basis choice of $M$ at the beginning, and the resulting quotient should be viewed as (the $\CC$-points of) the moduli space of 1-motives we just described. In fact, to get a fine moduli space, we need to quotient out the smaller group $P^{u}(\ZZ)$ first.

Now we present the main computation in \cite[Section 4]{BE20}, and we hope the illustrations above are helpful for the readers.
\begin{proposition}{\cite[Proposition 4.6]{BE20}}$\ $

    The quotient $P^{u}(\ZZ)\backslash D$ is the universal Poincar\'{e} torsor over $\HH_g$.
\end{proposition}
\begin{corollary}\label{unif}
The quotient $P^{u}(\ZZ)\M_{1,2g}(\RR)U(\CC)\backslash D$ is the universal abelian variety $\mathcal{A}$ over $\HH_g$, the quotient $P^{u}(\ZZ)\M_{2g,1}(\RR)U(\CC)\backslash D$ is the dual universal abelian variety $\mathcal{A}^{\vee}$ over $\HH_g$ and the quotient $P^{u}(\ZZ)U(\CC)\backslash D$ is the family $\mathcal{A}\times_{\HH_g}\mathcal{A}^{\vee}$ over $\HH_g$.
\end{corollary}
One can further quotient by some congruence subgroup of $\GSp(\Psi)(\QQ)$ to replace $\HH_g$ by some Siegel moduli space.

\subsection{Mixed Shimura Data of Ribet Sections}$\ $

In this subsection, we will describe the mixed Shimura data of the general Ribet sections we constructed before. To achieve this, we will first generalize the Hodge tensor \cite[(5.1.1)]{BE20} to the setting with the \textbf{Ribet property}(Notation \ref{ribet_prop}). Then we will explicitly compute the algebraic groups of the mixed Shimura data, which will enable us to prove the compatibility with the construction using algebraic geometry before.

Recall the bijection in Proposition \ref{D}. In this subsection, fix $\tilde{\tau}=(\tau,0,0,0)\in \HH_g\times \M_{1,g}(\CC)\times \M_{g,1}(\CC)\times \CC$.

\subsubsection{Description of a Hodge Tensor}$\ $

Recall the pure Hodge structure on $\ZZ^{2g}=\M_{2g,1}(\ZZ)$ corresponding to $\tau$ that we considered in Example \ref{ppav}. Then the abelian variety associated to this Hodge structure is $A_\tau=\M_{g,1}(\CC)/(1_g\ |-\tau)\M_{2g,1}(\ZZ)$. Identifying $\M_{2g,1}(\ZZ)^{\vee}$ with $\M_{1,2g}(\ZZ)$ via the symmetric bilinear pairing 
\[\M_{2g,1}(\ZZ)\times\M_{1,2g}(\ZZ)\rightarrow\ZZ,\ (x,y)\mapsto yx,\]
the induced Hodge structure on $\M_{1,2g}(\ZZ(1))$ corresponds to the dual of $A_\tau$. And the dual abelian variety is $A_\tau^{\vee}=\M_{1,g}(\CC)/\M_{1,2g}(\ZZ(1))\begin{pmatrix}\begin{smallmatrix}\tau\\1_g\end{smallmatrix}\end{pmatrix}$. Here we adapt the sign convention of \cite{BE20} in the definition of $A_\tau$ (note the minus sign).

Consider a \textbf{saturated} $\ZZ$-submodule $\Lambda$ of $\M_{2g,1}(\ZZ)\oplus\M_{1,2g}(\ZZ(1))$ of rank $2d$ satisfying two assumptions. The first is as follows.
\begin{assumption}\label{tensor}
Consider the injective map 
\begin{align*}
\ZZ(1)\oplus \Lambda\oplus \ZZ&\stackrel{\Theta_\Lambda}{\hookrightarrow} (\ZZ(1)\oplus\M_{2g,1}(\ZZ)\oplus\ZZ)
\oplus (\ZZ\oplus\M_{1,2g}(\ZZ(1))\oplus\ZZ(1))\\
(k,w_1,w_2,n)&\mapsto ((k,w_1,n),(-n,w_2,-k)).
\end{align*}
Define the weight filtration on $\ZZ(1)\oplus \Lambda\oplus \ZZ$ by 
\begin{align*}
&W_{-3}=\{0\},&W_{-2}=\ZZ(1),\\
&W_{-1}=\ZZ(1)\oplus\Lambda, &W_{0}=\ZZ(1)\oplus \Lambda\oplus \ZZ.
\end{align*}
The Hodge filtration induced by the above injection on $(\ZZ(1)\oplus \Lambda\oplus \ZZ)_{\CC}$ makes $\ZZ(1)\oplus \Lambda\oplus \ZZ$ a mixed Hodge structure.$\hfill\blacksquare$
\end{assumption}
\begin{remark}
    Recall that we have defined $M=\ZZ(1)\oplus\M_{2g,1}(\ZZ)\oplus\ZZ$ in Subsection \ref{MHS}. The summand $\ZZ\oplus\M_{1,2g}(\ZZ(1))\oplus\ZZ(1)$ in the definition of the map $\Theta_{\Lambda}$ should be viewed as $M^{\vee}(1)$.
\end{remark}
Under the above assumtion, $\Lambda$ is a sub-Hodge structure of $\M_{2g,1}(\ZZ)\oplus\M_{1,2g}(\ZZ(1))$. Since $\M_{2g,1}(\ZZ)\oplus\M_{1,2g}(\ZZ(1))$ is a polarizable Hodge structure, so is $\Lambda$. Therefore, $\Lambda$ defines a abelian subvariety $A'$ of $A_{\tau}\times A_{\tau}^{\vee}$. The second assumption on $\Lambda$ is as follows.
\begin{assumption}\label{ribetproperty}
    As a abelian subvariety of $A_{\tau}\times A_{\tau}^{\vee}$, $A'$ satisfies the \textbf{Ribet property}. $\hfill\blacksquare$
\end{assumption}

\begin{remark}
From the viewpoint of 1-motives, Assumption \ref{tensor} can be interpreted as follows. Despite the sign issue, when translated to 1-motives, the map $\Theta_{\Lambda}$ is meant to choose a point on the Poincar\'{e} torsor of $A_{\tau}$ whose projection to $A_{\tau}\times A_{\tau}^{\vee}$ lies in $A'$. Then under the uniformization map, $\tilde{\tau}$ should correspond to a point on the Ribet section $A'$. (In fact, the definition of $\tilde{\tau}$ is designed to make it a point over the origin of $A'$ and it is the unit element of the fibral $\gm$).
\end{remark}

Now we describe Assumption \ref{ribetproperty} using matrices. One can choose a basis of $\Lambda$ such that under the isomorphism with $\M_{2d,1}(\ZZ)$ there exists a $\tau_{\lambda}\in\delta_{\Lambda}^{-1}\HH_d$ with 
\[F^0\Lambda_{\CC}=\begin{pmatrix} \tau_{\Lambda}\\1_g
\end{pmatrix}\M_{d,1}(\CC).\] Here $\delta_{\Lambda}=\mathrm{diag}\{e_1,...,e_d\}$ for positive integers $e_1,...,e_d$ with $e_1|e_2|...|e_d$. And $\delta_{\Lambda}^{-1}\HH_g:=\{\tau\in\M_{g}(\CC):(\delta_{\Lambda}\tau)^{\top}=\delta_{\Lambda}\tau,\ \mathrm{Im}(\delta_{\Lambda}\tau)>0\}$. (See Remark \ref{poltype}.) Then there is a commutative diagram

\[\begin{tikzcd}
{\M_{2d,1}(\ZZ)} \arrow[d, "(1_g\ |-\tau_{\Lambda})"'] \arrow[r, "{(\alpha_{\ZZ},\beta_{\ZZ})}"] & {\M_{2g,1}(\ZZ)\oplus\M_{1,2g}(\ZZ(1))} \arrow[d, "(1_g\ |-\tau)"'] \arrow[d, "\begin{pmatrix}\begin{smallmatrix}
    \tau\\1_g
\end{smallmatrix}\end{pmatrix}"] \\
{\M_{d,1}(\CC)} \arrow[r, "{(\alpha_{\CC},\beta_{\CC})}"]                                          & {\M_{g,1}(\CC)\oplus\M_{1,g}(\CC)}                                                                                    
\end{tikzcd}\]
where the upper horizontal map is given by
\begin{align*}
    \M_{2d,1}(\ZZ)\rightarrow \M_{2g,1}(\ZZ),\ x\mapsto \alpha_\ZZ x
\end{align*}
and
\begin{align*}
    \M_{2d,1}(\ZZ)\rightarrow \M_{1,2g}(\ZZ(1)),\ y\mapsto 2\pi iy^\top\beta_\ZZ.
\end{align*}

These two maps correspond to the two projections $A'\stackrel{\alpha}{\rightarrow} A_{\tau}$ and $A'\stackrel{\beta}{\rightarrow} A_{\tau}^{\vee}$. Following exactly the same computations as in \cite[Proposition 4.8]{BE20}, one can determine the matrices corresponding to $\alpha^{\vee}$ and $\beta^{\vee}$, that is, the matrices in
\begin{align*}
    \M_{1,2g}(\ZZ(1))\rightarrow \M_{1,2d}(\ZZ(1)),\ x\mapsto x(\alpha^{\vee})_\ZZ 
\end{align*}
and
\begin{align*}
    \M_{2g,1}(\ZZ)\rightarrow \M_{1,2d}(\ZZ(1)),\ y\mapsto 2\pi iy^\top(\beta^{\vee})_\ZZ.
\end{align*}
\begin{lemma}\label{dual}
$(\alpha^{\vee})_{\ZZ}=\alpha_{\ZZ}$ and $(\beta^{\vee})_{\ZZ}=-\beta_{\ZZ}^{\top}$.
\end{lemma}


\subsubsection{Construction of the Mixed Shimura Data}$\ $\label{MSh}$ $

Recall that we have defined $M=\ZZ(1)\oplus\M_{2g,1}(\ZZ)\oplus\ZZ$  and an algebraic group $P\subset \GL(M)\times\GL(\ZZ(1))$ in Subsection \ref{MHS}.
\begin{notation}
    Let $Q_{\Lambda}$ be the subgroup scheme of $P$ fixing the map $\Theta_{\Lambda}$ in Assumption \ref{tensor}. Let $G_{\Lambda}=Q_{\Lambda}/(Q_{\Lambda}\cap P^u)\subset \GSp(\Psi)$ be the reductive quotient of $Q_{\Lambda}$.
\end{notation}
One can explicitly compute the functor of points of $Q_{\Lambda}$, which is crucial for the latter verifications. This computation is new. In \cite[Section 5]{BE20} a similar result is presented without further explanations.
\begin{lemma}\label{MT}
    For a $\ZZ[\frac{1}{2}]$-algebra $R$, the $R$-valued points of $Q_{\Lambda}$ is given by
    \[Q_{\Lambda}(R)=\left\{\begin{pmatrix}\mu(g)& v^\top\beta_\ZZ g&\frac{1}{2}v^\top\beta_\ZZ\alpha_\ZZ v\\
    0&g&\alpha_\ZZ v\\
    0&0&1\end{pmatrix}:v\in\M_{2d,1}(R), g\in G_{\Lambda}(R)\right\}.\]
\end{lemma}
\begin{proof}
For the reader's convenience, we repeat the definition of the map $\Theta_{\Lambda}$:
\begin{align*}
\ZZ(1)\oplus \Lambda\oplus \ZZ&\stackrel{\Theta_\Lambda}{\hookrightarrow} (\ZZ(1)\oplus\M_{2g,1}(\ZZ)\oplus\ZZ)
\oplus (\ZZ\oplus\M_{1,2g}(\ZZ(1))\oplus\ZZ(1))\\
(k,w_1,w_2,n)&\mapsto ((k, w_1,n),(-n,w_2,-k)).
\end{align*}
For
\[p=\begin{pmatrix}
    \mu(g)&x&z\\
    0&g&y\\
    0&0&1
\end{pmatrix}\in P(R),\]
\[p^{-1}=\begin{pmatrix}
    \mu(g)^{-1}&-\mu(g)^{-1}xg^{-1}&\mu(g)^{-1}(xg^{-1}y-z)\\
    0&g^{-1}&-g^{-1}y\\
    0&0&1
\end{pmatrix}.\]
The action of $p$ on $M_R$ is left multiplication by $p$ and the action of $p$ on $M^{\vee}(1)_R$ is right multiplication by $\mu(g)p^{-1}$.
The result is
\begin{itemize}
    \item $p\cdot(k, w_1,n)=(\mu(g)k+xw_1+nz,gw_1+ny,n)$;
    \item $p\cdot(-n,w_2,-k)=(-n,nxg^{-1}+\mu(g)w_2g^{-1},-n(xg^{-1}y-z)-\mu(g)w_2g^{-1}y-\mu(g)k)$.
\end{itemize}
Therefore, for $p\in Q_{\Lambda}(R)$, it is sufficient and necessary that for all $(k,w_1,w_2,n)\in (\ZZ(1)\oplus \Lambda\oplus \ZZ)_R$, the following two conditions
\begin{enumerate}
    \item $(gw_1+ny,nxg^{-1}+\mu(g)w_2g^{-1})\in\Lambda_R$;
    \item $xw_1+2nz=(nxg^{-1}+\mu(g)w_2g^{-1})y$.
    \end{enumerate}
hold.

For (1), since $(gw_1,\mu(g)w_2g^{-1})=g\cdot(w_1,w_2)\in\Lambda_{R}$, we have $(ny,nxg^{-1})\in\Lambda_{R}$. Hence $(y,xg^{-1})\in\Lambda_{R}$ since $\Lambda$ is saturated. Suppose $(y,xg^{-1})=(\alpha_{\ZZ}v,v^{\top}\beta_{\ZZ})$, $(gw_1,\mu(g)w_2g^{-1})=(\alpha_{\ZZ}u,u^{\top}\beta_{\ZZ})$. Then $xw_1=(xg^{-1})(gw_1)=v^{\top}\beta_{\ZZ}\alpha_{\ZZ}u$ and $(\mu(g)w_2g^{-1})y=u^{\top}\beta_{\ZZ}\alpha_{\ZZ}v$. The matrix $\beta_{\ZZ}\alpha_{\ZZ}$ is symmetric from $\alpha^{\vee}\beta+\beta^{\vee}\alpha=0$ and Lemma \ref{dual}. Hence $xw_1=(\mu(g)w_2g^{-1})y$ and from (2) we finally get $z=\frac{1}{2}xg^{-1}y=\frac{1}{2}v^{\top}\beta_{\ZZ}\alpha_{\ZZ}v$.
\end{proof}
Now we introduce the various mixed Shimura data involved.
\begin{notation}\label{groups}
$ $

\begin{itemize}
    \item The base $S$, which is a pure Shimura variety. 
    
    $\circ$ Define $\HH_{g,\Lambda}:=G_{\Lambda}(\RR)^{+}\cdot\tau\subset\HH_g$.
    
    $\circ$ The Shimura datum is thus $(G_{\Lambda},\HH_{g,\Lambda})$.

    \item The Poincar\'e torsor $(\Pt)_{S}$.

    $\circ$ Define $P_{\Lambda}$ to be the inverse image of $G_{\Lambda}$ under the quotient $P\rightarrow P/P^{u}=\GSp(\Psi)$, and $E_{\Lambda}:=P_{\Lambda}(\RR)^+\cdot U(\CC)\cdot\tilde{\tau}$.

    $\circ$ The mixed Shimura datum is thus $(P_\Lambda, E_\Lambda)$.

    \item The universal abelian scheme $(\Ag\times_{\ag} \Agv)_S$ over $S$. 

    $\circ$ Define $P_\Lambda':=(P_\Lambda \cdot U)/U$, and $E'_{\Lambda}:=P'_{\Lambda}(\RR)^+\cdot \tau'$. 

    $\circ$ The mixed Shimura datum is thus $(P_\Lambda', E'_{\Lambda})$.
    
    \item The Ribet section. 
    
    $\circ$ Define $D_{\Lambda}:=Q_{\Lambda}(\RR)^{+}\cdot\tilde{\tau}\subset D$.
    
    $\circ$ The mixed Shimura datum is thus $(Q_\Lambda, D_\Lambda)$.
    
    \item The image $\calA'$ of the Ribet section under $\pi:\Pt \to \Ag\times_{\ag} \Agv $ , which is an abelian subscheme of $(\Ag\times_{\ag} \Agv)_S$ over $S$. 
    
    $\circ$ Define $Q'_{\Lambda}:=(Q_{\Lambda}\cdot U) /U$, the image of $Q_{\Lambda}$ under the map $P\rightarrow P/U$. Also define $D'_{\Lambda}:=Q'_{\Lambda}(\RR)^{+}\cdot\tau’$, where $\tau'$ is the image of $\tau$ in $\HH_g\times \M_{1,g}(\CC)\times \M_{g,1}(\CC)$. 
    
    $\circ$ The mixed Shimura datum is thus $(Q'_\Lambda, D'_\Lambda)$.

    
\end{itemize}

\end{notation}
Finally, we have the comparison result.

\begin{proposition}
$(Q'_{\Lambda})^{u}(\ZZ)\backslash D’_{\Lambda}$ is an abelian subscheme of the abelian scheme $(P'_{\Lambda})^u(\ZZ)\backslash E'_{\Lambda}=(\Ag\times_{\HH_g}\Agv)_{\HH_{g,\Lambda}}$, and $(Q_{\Lambda})^{u}(\ZZ)\backslash D_{\Lambda}$ is the Ribet section over it as constructed in Section \ref{section2}.
\end{proposition}

\begin{proof}
Let $(\tilde{Q}_{\Lambda})^{u}(2\ZZ)$ be the group
\[\left\{\begin{pmatrix}1& v^\top\beta_\ZZ&\frac{1}{2}v^\top\beta_\ZZ\alpha_\ZZ v\\
    0&1&\alpha_\ZZ v\\
    0&0&1\end{pmatrix}:v\in\M_{2d,1}(2\ZZ)\right\}\]
and let $(\tilde{Q}'_{\Lambda})^{u}(2\ZZ)$ be it is image in $(P/U)(\ZZ)$.
Then $(\tilde{Q}_{\Lambda})^{u}(2\ZZ)$ is a subgroup of both $(P_{\Lambda})^{u}(\ZZ)$ and $(Q'_{\Lambda})^u(\ZZ)$. We summarize the above constructions in the following diagram.
\[
\begin{tikzcd}
(\tilde{Q}_{\Lambda})^{u}(2\ZZ)\backslash D_{\Lambda} \arrow[d, "\cong"'] \arrow[r] & (Q_{\Lambda})^{u}(\ZZ)\backslash D_{\Lambda} \arrow[d] \arrow[r] & (P_{\Lambda})^u(\ZZ)\backslash E_{\Lambda} \arrow[d] \\
(\tilde{Q}'_{\Lambda})^{u}(2\ZZ)\backslash D’_{\Lambda} \arrow[r, "{[2]}"]          & (Q'_{\Lambda})^{u}(\ZZ)\backslash D’_{\Lambda} \arrow[r]         & (P'_{\Lambda})^u(\ZZ)\backslash E'_{\Lambda}    
\end{tikzcd}\]
Compare the above diagram with the diagram in Section \ref{section2}, which reads now:
\[
\begin{tikzcd}
{[2]_{A'}^*(\pt|_{\calA'})} \arrow[d] \arrow[r, "\pi"] \arrow[dr, phantom,"\ulcorner",very near start]            & \pt|_{\calA'} \arrow[d] \arrow[r] \arrow[dr, phantom,"\ulcorner",very near start]& \pt \arrow[d]    \\
\calA' \arrow[r, "{[2]_{\calA'/S}}"] \arrow[u, "\tilde{r}_{\calA'}", dotted, bend left] & \calA' \arrow[r, hook]              & (\Ag\times_{\ag} \Agv)_S
\end{tikzcd}
\]
we see that $(Q_{\Lambda})^{u}(\ZZ)\backslash D_{\Lambda}$ is up to a global invertible function on $\HH_{g,\Lambda}$ equal to the Ribet section we constructed using algebraic geometry. Due to the choice of $\tilde{\tau}$ at the beginning, this function is identically 1 since the value of $(Q_{\Lambda})^{u}(\ZZ)\backslash D_{\Lambda}$ above the identity section of $(Q'_{\Lambda})^{u}(\ZZ)\backslash D’_{\Lambda}$ is 1.
\end{proof}

\section{Classification of Special Subvarieties of the Poincar\'{e} Torsor}\label{section4}

We shall finish the proof of Theorem \ref{thm_Pt} in this section. Let us recall that, there are projections $\Pi: \Pt\rightarrow \ag$, $\pi: \Pt\rightarrow \Ag\times_{\ag}\Agv$ and $\pr: \Ag\times_{\ag}\Agv\rightarrow \ag$.

Let $\Pi(T)=S$ a special subvariety of the base $\ag$, and $\pi(T)=W$ a special subvariety of $\Ag\times_{\ag} \Agv$. We thus restrict everything to $S$. To make notations compatible with the previous section, we denote the pure Shimura datum of $S$ by $(G_\Lambda, \HH_{g,\Lambda})$.

An analogous result of \cite[Proposition 3.4]{gaoziyang} says that, $W=\calA'+t$ is a torsion translate of some abelian subscheme $\mathcal A'\hookrightarrow \Ag\times_{\ag}\Agv$ over $S$. 

The proof of the classification is thus divided into three steps. 
\begin{itemize}

\item Firstly we construct $T'$, a torsion translate of $T$, which is still a special subvariety of $\Pl$ whose image $\calA'=\pi(T)$ is an abelian subscheme of $\Ag\times_{\ag}\Agv$ over $S$. From that, we may assume $T'=T$ and $\calA'=W$. 

\item Secondly, for the morphism of mixed Shimura datum $(Q,X)\to (Q',X')$ corresponds to $\pi|_{T}:T\to \calA'$, there is a commutative diagram of algebraic group homomorphisms,
\begin{center}

\begin{tikzcd}
Q \arrow[r] \arrow[d, hook] & Q' \arrow[d, hook] \\
P \arrow[r]                   & P/U                
\end{tikzcd}
\end{center}
where the horizontal arrows are embeddings of subgroups. Recall that $U\cong\mathbb G_a$, then $K=\ker(Q\to Q')$ is either $0$ or $\mathbb G_a$. 

\item Finally, the proof is finished by discussing these two possibilities of $K$.

\end{itemize}

\subsection{Moving by Torsion Sections}\label{moving}$ $

In this section, we review the known results of torsion multi-sections and explain the details in the first step of the proof.

\begin{definition}
    Let $\calA/S$ be a commutative group scheme. A torsion section is a torsion element in $\sigma\in \calA(S)$.

    If $\calA/S$ is an abelian variety, a torsion multi-section is a finite surjective map $S'\to S$ with a section $\sigma \in \calA(S')$.
\end{definition}

We now give the result on classification of Shimura subdatum $(Q',X')$ (corresponding to $t+\calA'$) of $(P',E')$ (corresponding to $\Ag\times_\ag \Agv$). The result is essentially \cite[Proposition 3.4]{gaoziyang}.
\begin{proposition}
To give a mixed Shimura subdatum $(Q',X')$ of $(P',E')$ is equivalent to giving:
\begin{itemize}
    \item a pure Shimura subdatum $(G_\Lambda,\HH_{g,\Lambda})$ of $(\GSp_{\Psi}, \HH_{g})$;
    \item a $Q'$-submodule $V_{Q'}$ of $\M_{1,2g}\times \M_{2g,1}$;
    \item a vector $(\bar v,\bar {v^\vee})\in ((\M_{1,2g}\times \M_{2g,1})/V_{Q'} )(\QQ)$.
\end{itemize}
In such case, the embedding $Q'\to P'$ is the conjugation of the usual one by 
    \[\begin{pmatrix}1&v^\vee &\frac{1}{2}v^\vee\cdot v\\
    0&1& v\\
    0&0&1 \end{pmatrix},\]
where $v\in \M_{1,2g}(\QQ)$ is any lifting of $\bar v$, and similar for $v^\vee$. 

\end{proposition}

\begin{remark}
    Corrspondingly, any mixed Shimura subvariety of $\Ag\times_\ag \Agv$ is determined by its Shimura datum $(Q',X')$ as follows: 
    \begin{itemize}
    \item the pure Shimura subdatum $(G_\Lambda,\HH_{g,\Lambda})$ of $(\GSp_{\Psi}, \HH_{g})$ gives the base $S$, which is a special subvariety of $\ag$;
    \item the $Q'$-submodule $V_{Q'}$ of $\M_{1,2g}\times \M_{2g,1}$ gives an abelian subscheme $\calA'$ of $(\Ag\times_\ag \Agv)_S$.
    \item a vector $(\bar v,\bar {v^\vee})\in ((\M_{1,2g}\times \M_{2g,1})/V_{Q'} )(\QQ)$, determines the torsion multi-section $t$.
\end{itemize}
    Thus the torsion multi-section $t$ is given by a section after a base change by a finite surjective Shimura morphism $S'\to S$. Moreover, $S'$ and $S$ have the same Shimura datum (so they differ by level structure).
    
    Finally, the mixed Shimura subvariety, denoted by $\calA'+t$, will be the image of the following  composition: $$t+\calA'_{S'}\hookrightarrow(\Ag\times_\ag \Agv)_{S'}\to (\Ag\times_\ag \Agv)_{S}$$
\end{remark}

For simplicity, one makes a base change by $S'\to S$ to assume $S'=S$, and $t\in (\Ag\times_\ag \Agv)(S)$ a single-valued section (this amounts to get into a deeper level structure). In the following, $\Ag$, $\Agv$ and $\Pt$ would be understood as $(\Ag)_S$, $(\Agv)_S$, and $(\Pt)_S$, respectively.

Now comes the moving step. Write $t=(x,y)$, where $x\in \Ag(S)$, $y\in \Agv(S)$ are torsion elements. Let $ny=0$, and $y$ corresponds to a torsion line bundle $\lb$ on $\Ag$, with $\lb ^{\otimes n}=0$. Then $[n]^*\lb\cong \calO_\Ag$. Using the rigidification of $\Pl$, one obtain a section $\iota:\Ag\times_S y \to [n]^*(\Pt)_y$, splitting the first arrow:

\begin{center}
\begin{tikzcd}
1 \arrow[r] & \Gm \arrow[r] \arrow[d] & {[n]^*(\Pt)_y} \arrow[r] \arrow[d] \arrow[dr, phantom,"\ulcorner",very near start]& \Ag\times_S y \arrow[r] \arrow[d, "{[n]}"] \arrow[l, "\iota"', dotted, bend right] & 1 \\
1 \arrow[r] & \Gm \arrow[r]           & (\Pt)_y \arrow[r]           & \Ag\times_S y \arrow[r]                                                            & 1
\end{tikzcd}
\end{center}

One lifts $\Ag$ to a closed subscheme $\tilde y$ of $[n]^*\Pt$, by taking the image of the closed embedding:$$\Ag=\Ag\times_S y\stackrel\iota\to[n]^*(\Pt)_y\hookrightarrow[n]^*(\Pt).$$
The projection of $\tilde y$ to $\Agv$ is $y$. 

Since $[n]^*\Pt$ is a group scheme over $\Ag$, there is a subtraction map:$$\mathrm{minus}:[n]^*\Pt\times_\Ag [n]^*\Pt\stackrel {(\id,[-1])}\longrightarrow[n]^*\Pt\times_\Ag[n]^*\Pt\stackrel m\to[n]^*\Pt.$$ 
One considers $T_1$, which is defined to be the image of $$T'\times_\Ag\tilde y \hookrightarrow [n]^*\Pt\times_\Ag [n]^*\Pt\stackrel {\mathrm{minus}}\longrightarrow[n]^*\Pt\to\Pt,$$
where $T'$ is an irreducible component of the inverse image of $T$ under the map $[n]^*\Pt\to\Pt$.
Since $\tilde y$ is isomorphic to $\Ag$, and $T'\times_\Ag \tilde y$ is isomorphic to $T'$, we conclude that the image $T_1$ is irreducible.

Now by construction, the morphisms involved are mixed Shimura morphisms. Thus $T_1$ is a mixed Shimura subvariety of $\Pt$. We now compute the image of $T_1$ under the projection $\Pt\stackrel \pi \to  \Ag\times_S\Agv$. In the following diagram, we set $\calB:=\Ag\times_S\Agv$ for the brevity of notations.

\[\begin{tikzcd}[row sep=5ex, column sep=3ex]
T'\times_\Ag \tilde y \arrow[rd, hook] \arrow[rrr, maps to] \arrow[ddd, maps to] &                                                                               &                                & T_1 \arrow[d, hook] \arrow[ddd, maps to, bend left] \\
                                                                                & {[n]^*\Pt\times_\Ag[n]^* \Pt} \arrow[r, "\mathrm{minus}"] \arrow[d] \arrow[d] & {[n]^*\Pt} \arrow[d] \arrow[r] & \Pt \arrow[d]                                       \\
                                                                                & \calB\times_\Ag\calB \arrow[r, "\mathrm{minus}"]                              & \calB \arrow[r, "{(n,\id_{\Agv})}"]   & \calB                                               \\
{(\calA'+(x,y))\times_\Ag(\Ag\times_S y)} \arrow[ru, hook] \arrow[rrr, maps to] &                                                                               &                                & {\calA'+(x,0)} \arrow[u, hook]                     
\end{tikzcd}
\]
From the commutativity of the outer square of the diagram, we conclude that the image of $T_1$ is exactly $\calA'+(x,0)$. Apply similar operations for $x$. Thus we may assume that $\calA'=\pi(T)$ is an abelian subscheme.

\begin{remark}\label{trans}
\begin{enumerate}
    \item For the purpose of classifying special subvarieties, the above translation process is essentially invertible. In fact, one can recover the original $T$ as follows. First, take the inverse image of $T_1$ under the \textit{finite} morphism $[n]^*\Pt\rightarrow\Pt$. This may be reducible, but the dimension of it is the same as $T$. Due to dimensional reasons, one of its irreducible components must be the image of the map
    \[T'\times_\Ag\tilde y \hookrightarrow [n]^*\Pt\times_\Ag [n]^*\Pt\stackrel {\mathrm{minus}}\longrightarrow[n]^*\Pt.\]
    Translate it back and one recovers $T’$. Then the image of $T'$ under the map $[n]^*\Pt\to\Pt$ is $T$ again due to dimensional reasons. 

    \item We now summarize the meaning of the term “torsion translate”. In the above constructions, we first choose an irreducible component of the inverse image of $T$ under the map $[n]^*\Pt\to\Pt$. Translate it by a section of $\Ag\times_S y$ and then project down to $\Pt$ to get $T_1$, a torsion translate of $T$. It is easy to formulate a general situation in exactly the same patterns as above. But we will not do this as this will introduce plenty of notations.

    \item Careful readers may be concerned with the \textit{order} of the two translation processes. Translate by $-y$ first or translate by $-x$ first will only differ by a torsion multi-section in the $\gm$’s direction, which is the Weil pairing of $-x$ and $-y$. Such a difference actually appears in the proofs in Subsection \ref{subsection4.3}, and we will tackle this issue there.
\end{enumerate}

\end{remark}

\subsection{The case of $K=\Ga$}$\ $

We assume in this and the next subsection that the image of the special subvariety $T$ under the map $\pi: \Pt\rightarrow \Ag\times_{\ag}\Agv$ is an abelian subscheme.

In the case of the title of this subsection, the result is nearly trivial. The special subvariety we seed will just be an inverse image of an abelian subscheme of $\Ag\times_{\ag}\Agv$ under the projection $\pi: \Pt\rightarrow \Ag\times_{\ag}\Agv$, whose projection to $\ag$ is a pure Shimura subvariety of $\ag$.

\subsection{The case of $K=0$}\label{subsection4.3}
$\ $

This is the essential case of the entire proof. Let $S$ be the projection of $T$ to $\ag$. To simplify notations, in this subsection, we will denote $(\Ag)_S$, $(\Agv)_S$, and $(\Pt)_S$ by $\Ag$, $\Agv$ and $\Pt$, respectively.

Let $\calA'$ be the projection of $T$ to $\Ag\times_{S}\Agv$. To determine the abelian subscheme $A'$ of $\Ag\times_{S}\Agv$, the following is the crucial observation.
\begin{proposition}\label{char}
    In the notations as above, $\calA'$ is an abelian subscheme of $\Ag\times_{S}\Agv$ satisfying the \textbf{Ribet property} (Proposition-Definition \ref{ribet_prop}).
\end{proposition}
\begin{proof}
From \cite[Facts 2.6(a)]{pink04} one deduces that $T\rightarrow \calA'$ is a \textit{finite} morphism. Consider the following cartesian diagram.
\[
\begin{tikzcd}
\Pt|_{\calA'}\times_{\calA'}T \arrow[d] \arrow[r] \arrow[dr, phantom,"\ulcorner",very near start]       & \Pt|_{\calA'} \arrow[d] \\
T \arrow[r, "f"] \arrow[u, "\iota_T", dotted, bend left] & \calA'                 
\end{tikzcd}\]
Here $\iota_T$ is the tautological section. Note that the existence of such a section to the $\gm$-torsor $\Pt|_{\calA'}\times_{\calA'}T$ forces it to be a trivial torsor. Hence the pullback line bundle $f^*(\Pl|_{\calA'})$ is trivial. For every geometric point $x$ of $S$, the line buncle $(f^*(\Pl|_{\calA'}))_x\cong f_x^*(\Pl|_{\calA'_x})$ is thus trivial. Since $f_x:T_x\rightarrow \calA'_x$ is finite surjective, by the projection formula \cite[Proposition 2.5(c)]{fulton}, one computes
\begin{align*}
\deg(f_x)\cdot c_1(\Pl|_{\calA'_x})\cap [\calA'_x]&=c_1(\Pl|_{\calA'_x})\cap f_{x,*}[T_x]\\
&=f_{x,*}(f_x^*(c_1(f_x^*(\Pl|_{\calA'_x})))\cap[T_x])\\
&=0.
\end{align*}
Since N\'eron-Severi group of any abelian variety is a torsion free $\ZZ$-module, $\Pl|_{\calA'_x}$ lies in $\pic^0_{\calA'_{x}/x}(\kappa(x))$ for every geometric point of $S$. Thus by the moduli interpretation of $\pic^0_{\calA'/S}$ one deduces that $\Pl|_{\calA'}$ lie in $\pic^0_{\calA'/S}(S)$. Hence $\calA'$ satisfies the Ribet property.
\end{proof}

Recall from Notation \ref{groups} that the mixed Shimura datum of $\calA'$ is of the form $Q'_{\Lambda}$ for some $\Lambda$ considered in Subsection \ref{MSh}. Let $P_T$ be the algebraic group in the mixed Shimura datum of $T$. Then $T\hookrightarrow \Pt$ corresponds to an inclusion of $\QQ$-algebraic groups $P_T\hookrightarrow P$ such that the composite
\[P_T\hookrightarrow P\twoheadrightarrow P/U\]
is an isomorphism onto its image $Q'_{\Lambda}$. Here is a lemma describing the $R$-valued points of $P_{T}$ for any $\QQ$-algebra $R$.
\begin{lemma}\label{fcneq}
    For any $\QQ$-algebra $R$,
    \[P_{T}(R)=\left\{\begin{pmatrix}\mu(g)& v^\top\beta_\ZZ g&\frac{1}{2}v^\top\beta_\ZZ\alpha_\ZZ v+(\mu(g)-1)z\\
    0&g&\alpha_\ZZ v\\
    0&0&1\end{pmatrix}:v\in\M_{2d,1}(R), g\in G_{\Lambda}(R)\right\}\]
for some $z\in \QQ$. 
\end{lemma}
\begin{proof}
 For any $\QQ$-algebra $R$, suppose 
 \[P_{T}(R)=\left\{\begin{pmatrix}\mu(g)& v^\top\beta_\ZZ g&f(v,g) \\
    0&g&\alpha_\ZZ v\\
    0&0&1\end{pmatrix}:v\in\M_{2d,1}(R), g\in G_{\Lambda}(R)\right\}.\]Thus $f:G_\Lambda\times_\QQ\M_{2d,1}\rightarrow\mathbb{A}^1_\QQ$ is a morphism of schemes over $\QQ$. In fact, one defines $f$ by the composition:
    $$P/U\cong P_T\hookrightarrow \GL_{2n+2}\to\mathbb A^1,$$
    where the last morphism is the projection to the entry at the upper right corner.
    The goal is thus to determine the precise form of $f$, which determines the group $P_T$. Using the group law of $P_T$, for $v,w\in\M_{2d,1}(R), g,h\in G_{\Lambda}(R)$ we have an equation
    \[\mu(g)f(w,h)+v^\top\beta_\ZZ g \alpha_\ZZ w+f(v,g)=f(x,gh),\]
    where $x$ is the unique element in $\M_{2d,1}(R)$ satisfying the following two equations
    \begin{itemize}
        \item $\alpha_\ZZ x=g\alpha_\ZZ w+\alpha_\ZZ v$;
        \item $x^\top\beta_{\ZZ}=\mu(g)w^\top\beta_\ZZ g^{-1}+v^\top\beta_\ZZ$.
    \end{itemize}
    By plugging into values, the equation above implies the following three specific equations
    \begin{enumerate}
        \item $\mu(g)f(0,h)+f(0,g)=f(0,gh)$;
        \item $f(0,h)+f(v,1)=f(v,h)$;
        \item $f(w,1)+v^\top \beta_\ZZ\alpha_\ZZ w+f(v,1)=f(w+v,1)$.
    \end{enumerate}
    Actually, one also checks these three equations together to imply the original one. We leave the explicit calculations for the readers.

    From the definition of $G_{\Lambda}$, it is easy to see that $G_{\Lambda}$ contains all scalar matrices. Plugging $g=2\cdot \id$ into equation (1), we obtain $$
    4f(0,h)+f(0,2\cdot\id)=f(0,2\cdot h)=f(0,h\cdot 2)=\mu(h)f(0,2\cdot\id)+f(0,h).
    $$
    Thus, we can take $z:=f(0,2\cdot \id)/3\in \QQ$, such that $$f(0,h)=\frac {f(0,2\cdot \id)} 3 \cdot (\mu(h)-1)=(\mu(h)-1)z.$$

    We set $\varphi(v):=f(v,1)-\frac 1 2 v^\top \beta_\ZZ\alpha_\ZZ v$. Equation (3) is equivalent to$$\varphi(v)+\varphi(w)=\varphi(v+w),$$ because $\beta_\ZZ\alpha_\ZZ$ is a symmetric matrix. From this we conclude there is an $m\in M_{1,2d}(\QQ)$, such that $mv=\varphi(v)$.

    Combining these two results with equation (2), we conclude that $$
    f(v,g)=f(v,1)+f(0,g)=\frac 1 2 v^\top \beta_\ZZ\alpha_\ZZ v +mv+(\mu(g)-1)z.
    $$


    Now it suffices to prove that $m=0$. Without loss of generality, we may assume that $z=0$. Then up to a base change by an isogeny of $\mathcal{A}'$, the two subvarieties of $\Pt$ through the point corresponding to $\tilde{\tau}$ defined by the groups $P_T$ and $Q_{\Lambda}$ respectively are two sections of $\Pt|_{\mathcal{A}'}$ over $\mathcal{A}'$. Trivialize the Poincar\'{e} torsor by the Ribet section (the section corresponding to the group $Q_{\Lambda}$). The section defined by the group $P_T$ thus defines a morphism from $\mathcal{A}'$ to $(\gm)_{\mathcal{A}'}$. Choose a $\CC$-point of the base Shimura variety and restrict everything to the fiber over it. Then such a morphism must be constant, which is impossible unless $m=0$.
\end{proof}

In geometric terms, this lemma and the above discussion amount to saying that $T$ is the translation of the Ribet section over $\calA'$ by an element in ${\gm}_{\mathrm{tors}}$ by the natural $\gm$-action on the Poincar\'{e} torsor. The classification in Theorem \ref{thm_Pt} is thus finished at this stage.

\section{Classifications in the General Case}\label{section5}
Fix a positive integer $n$. In the following, $(X)^{\times_{Y}n}$ means the $n$-th power self-product $\underbrace{X\times_{Y}X\times_{Y}\cdots \times_{Y}X}_{n \ \mathrm{times}}$. Recall the definition of the universal semi-abelian variety $\pgn$. It is defined by the following Cartesian diagram. 

\[\begin{tikzcd}
\pgn \arrow[rr] \arrow[d] \arrow[drr, phantom,"\ulcorner",very near start]                                              &  & (\Pt)^{\times_{S_g}n} \arrow[d]       \\
\Ag\times_{S_g}(\Agv)^{[n]} \arrow[rr, "{(\Delta_{\Ag},\mathrm{id})}"] &  & (\Ag\times_{S_g}\Agv)^{\times_{S_g}n}
\end{tikzcd}\]

The classification of special subvarieties of $\pgn$ roughly follows the same pattern as in previous sections. We first describe the situation in Subsection \ref{generalMHS} using mixed Hodge structures or in other words, describe the various mixed Shimura data involved. Then in Subsection \ref{isogeny}, \ref{classify} we use an isogeny trick to reduce the classification to the case of the Poincar\'{e} torsor.

\subsection{Mixed Hodge structures}\label{generalMHS}$\ $

In this subsection, we generalize some constructions in subsection \ref{MHS} to the case of $\pgn$. Since the proofs are exactly the same, we will simply list the generalized results.

Let $g\geq 1$ be an integer and $M=\ZZ(1)^n\oplus \ZZ^{2g}\oplus \ZZ$ with standard basis $2\pi i e_0,...,2\pi ie_{n-1},e_n,...,e_{n+2g}$. Impose a weight filtration on $M$ as follows:
\begin{align*}
W_{-3}M&=\{0\},\\
W_{-2}M&=\ZZ2\pi i e_0\oplus\cdots \oplus \ZZ2\pi i e_{n-1},\\
W_{-1}M&=\ZZ2\pi i e_0\oplus\cdots \oplus \ZZ2\pi i e_{n-1}\oplus \ZZ e_n\oplus\cdots\oplus \ZZ e_{n+2g-1}, \\
W_{0}M&=M.
\end{align*}
Let $D$ be the set of decreasing filtrations $F^{\bullet}$ on $M_{\CC}$ such that $(M,W_{\bullet},F^{\bullet})$ is a $\ZZ$-mixed Hodge structure of type $\{(-1,-1),(-1,0),(0,-1),(0,0)\}$ and such that 
\[\Psi:\ZZ^{2g}\otimes \ZZ^{2g}\rightarrow \ZZ(1),\ x\otimes y\mapsto 2\pi i x^\top\begin{pmatrix}\begin{smallmatrix} 0&-1_g\\1_g&0
\end{smallmatrix}\end{pmatrix}y\]is a polarization of pure Hodge structure on $\Gr^{W}_{-1}M$ under the isomorphism $\Gr^{W}_{-1}M\cong \ZZ^{2g}$ given by the chosen basis.

\begin{lemma}$\ $

There is a bijection $\HH_g\times \M_{1,g}(\CC)^n\times \M_{g,1}(\CC)\times \CC^n \mapsto D$,
\[(\tau,\underline{u},v,\underline{w})\mapsto F^0 M_{\CC}=\begin{pmatrix}
    \underline{u}&\underline{w}\\
    \tau& v\\
    1_g&0\\
    0&1
\end{pmatrix}\CC^{g+1}.\]
\end{lemma}

Let $P$ be the subgroup scheme of $\GL(M)\times \GL(\ZZ(1))$ that fixes
\begin{itemize}
    \item $W_{\bullet}$, 
    \item $\ZZ(1)^n\rightarrow W_{-2}M$, $(2\pi i a_j)_{0\leq j\leq n-1}\mapsto (2\pi i a_je_j)_{0\leq j\leq n-1}$, 
    \item $\ZZ\rightarrow \Gr_0^{W}M$, $a\mapsto ae_{n+2g}$,
    \item $\Psi:\Gr^{W}_{-1}M\otimes \Gr^{W}_{-1}M\rightarrow \ZZ(1)$. 
\end{itemize}
\begin{lemma}
    The $R$-valued points of $P$ for a ring $R$ are
\begin{align*}
    P(R)=\left\{\begin{pmatrix}
    \mu(g)1_n&\underline{x}&\underline{z}\\
    0&g&y\\
    0&0&1
\end{pmatrix}:
    (g,\mu(g))\in\GSp(\Psi)(R),
\underline{x}\in\M_{n,2g}(R),y\in\M_{2g,1}(R),\underline{z}\in R^n 
 \right\}.
\end{align*}
\end{lemma}
The $R$-valued points of its unipotent radical are
\begin{align*}
    P^u(R)=\left\{\begin{pmatrix}
    1_n&\underline{x}&\underline{z}\\
    0&1&y\\
    0&0&1
\end{pmatrix}:
\underline{x}\in\M_{n,2g}(R),y\in\M_{2g,1}(R),\underline{z}\in R^n 
 \right\}.
\end{align*}
And $P$ has a normal subgroup $U$, whose $R$-valued points are
\begin{align*}
    U(R)=\left\{\begin{pmatrix}
    1_n&0&\underline{z}\\
    0&1&0\\
    0&0&1
\end{pmatrix}:\underline{z}\in R^n 
 \right\}.
\end{align*}
\begin{proposition}$\ $

    The quotient $P^{u}(\ZZ)\backslash D$ is $(\pgn)_{\HH_g}$.
\end{proposition}
\begin{corollary}
The quotient $P^{u}(\ZZ)\M_{n,2g}(\RR)U(\CC)\backslash D$ is the universal abelian variety $\mathcal{A}$ over $\HH_g$, the quotient $P^{u}(\ZZ)\M_{2g,1}(\RR)U(\CC)\backslash D$ is $n$-th fold fiber product of the dual universal abelian variety $\mathcal{A}^{\vee}$ over $\HH_g$ and the quotient $P^{u}(\ZZ)U(\CC)\backslash D$ is the family $\mathcal{A}\times_{\HH_g}(\mathcal{A}^{\vee})^{\times_{\mathbb{H}_g}n}$ over $\HH_g$.
\end{corollary}
\subsection{An Isogeny Trick}\label{isogeny}$ $

Now let $T$ be a special subvariety of $\pgn$ with mixed Shimura datum $(P_T, D_T)$. Denote $U'$ to be the subgroup $P_T\cap U_\QQ$ of $U_\QQ$. One is led to the idea of “quotient out the subtorus of $\pgn$ corresponding to $U'$”. However, since the position of $U'$ in $U$ may be arbitrary, in general, there is no such subgroup, especially in the family situation.

To settle this, we will consider an outer automorphism of the group $P_\QQ$ to transform $U'$ to a “standard coordinate space” by a Shimura morphism which is geometrically just an isogeny of the group scheme $\pgn\cong \underbrace{\Pt\times_{\Ag}\Pt\times_{\Ag}\cdots \times_{\Ag}\Pt}_{n \ \mathrm{times}}$. 

Choose a matrix $h\in\GL_n(\QQ)\cap \M_{n,n}(\ZZ)$ such that $hU'=\mathrm{im}(\Ga^{k}\hookrightarrow \Ga^n\cong U)$, where $k=\dim(U')$ and the inclusion is the standard inclusion to the first $k$ coordinates. Consider the automorphism $P_\QQ\stackrel{\mathrm{int}(h)}{\rightarrow}P_{\QQ}$ which on $R$-points for an $\QQ$-algebra $R$ is given by
\[\begin{pmatrix}
    \mu(g)1_n&\underline{x}&\underline{z}\\
    0&g&y\\
    0&0&1
\end{pmatrix}\rightarrow \begin{pmatrix}
    \mu(g)1_n&h\underline{x}&h\underline{z}\\
    0&g&y\\
    0&0&1
\end{pmatrix}=
\begin{pmatrix}
    h&0&0\\
    0&1_{2g}&0\\
    0&0&1
\end{pmatrix}
\begin{pmatrix}
    \mu(g)1_n&\underline{x}&\underline{z}\\
    0&g&y\\
    0&0&1
\end{pmatrix}\begin{pmatrix}
    h&0&0\\
    0&1_{2g}&0\\
    0&0&1
\end{pmatrix}^{-1}.\]
Note that although this automorphism is written as if it is an inner automorphism, it is actually an outer automorphism of $P_\QQ$. Geometrically, this corresponds to the isogeny 
\[
\begin{tikzcd}
\underbrace{\Pt\times_{\Ag}\Pt\times_{\Ag}\cdots \times_{\Ag}\Pt}_{n \ \mathrm{times}} \arrow[rd] \arrow[rr, "{[h]}"] &     & \underbrace{\Pt\times_{\Ag}\Pt\times_{\Ag}\cdots \times_{\Ag}\Pt}_{n \ \mathrm{times}} \arrow[ld] \\
                                                     & \Ag &                                 
\end{tikzcd}.\]

Replacing $T$ by $[h](T)$, we may assume that $U'\subset U$ is the coordinate space of the first $k$ coordinates. Now consider the projection 
\[\underbrace{\Pt\times_{\Ag}\Pt\times_{\Ag}\cdots \times_{\Ag}\Pt}_{n \ \mathrm{times}}\longrightarrow (\Ag\times_{\ag}(\Agv)^{[k]})\times_{\Ag} \underbrace{\Pt\times_{\Ag}\Pt\times_{\Ag}\cdots \times_{\Ag}\Pt}_{n-k \ \mathrm{times}}.\]
The special subvariety $T$ will be the inverse image of a special subvariety $\tilde{T}$ of \[(\Ag\times_{\ag}(\Agv)^{[k]})\times_{\Ag} \underbrace{\Pt\times_{\Ag}\Pt\times_{\Ag}\cdots \times_{\Ag}\Pt}_{n-k \ \mathrm{times}}.\] The goal is thus to classify such $\tilde{T}$.

\subsection{Finish the Classification}\label{classify}$ $

First consider the further projection $T'$ of $T$ to $\underbrace{\Pt\times_{\Ag}\Pt\times_{\Ag}\cdots \times_{\Ag}\Pt}_{n-k \ \mathrm{times}}$. This is a subvariety of $\underbrace{\Pt\times_{\Ag}\Pt\times_{\Ag}\cdots \times_{\Ag}\Pt}_{n-k \ \mathrm{times}}$ finite over $\Ag\times_{S_g} (\Agv)^{[n-k]}$. Thus if the image of $T'$ in $\Ag\times_{S_g} (\Agv)^{[n-k]}$, denoted by $\mathcal{A}'$ is an abelian subscheme of $\Ag\times_{S_g} (\Agv)^{[n-k]}$, then the image of $\mathcal{A}'$ in each $\Ag\times_{S_g} \Agv$ will be an abelian subscheme satisfying the Ribet property. A computation completely analogous to Lemma \ref{fcneq} and a discussion like the one after Lemma \ref{fcneq} will reveal that $T'$ is a translation of the fiber product over $\Ag$ of the various Ribet sections by a torsion element in the $\Gm$-direction. In general, if $\mathcal{A}'$ is just a torsion translate of an abelian subscheme, one can still easily argue like Subsection \ref{moving} that $T'$ is the fiber product over $\Ag$ of several special subvarieties of $\Pt$.

Considering $\tilde{T}$, the image of it in $(\Ag\times_{\ag}(\Agv)^{[k]})$ will be a torsion translate of an abelian subscheme, which we will denote by $\tilde{\mathcal{A}}$. Counting dimensions, one can easily get that $\tilde{T}$ is just an irreducible component of $\tilde{\mathcal{A}}\times_{\Ag} T'$. 

Combining the above discussions, $T$ is just an irreducible component of $T_1\times_{\Ag}T_2\times_{\Ag}\cdots\times_{\Ag}T_n$, where the $T_i$'s are special subvarieties of $\Pt$ whose images in $\Ag$ are the same. Thus the Theorem \ref{general} is proved. 
 
\begin{remark}
To make the discussions above clearer, we present the typical form (up to conjugation) of the group $P_T$ below. Let $d\leq(n+1)g$ be a positive integer, $\alpha\in \M_{2g,2d}(\ZZ)$, $\beta_i\in\M_{2d,2g}(\ZZ), 1\leq i\leq n$ be matrices such that for $k+1\leq i\leq n$, $\beta_i\alpha$ is symmetric. Then consider the $\QQ$-group $Q$ whose $R$-valued points for an $\QQ$-algebra $R$ are given by
\begin{align*}
    Q(R)=\left\{\begin{pmatrix}
    \mu(g)1_k&0&(v^\top\beta_i g)_{1\leq i\leq k}&\underline{z}\\
    0&\mu(g)1_{n-k}&(v^\top\beta_i g)_{k+1\leq i\leq n}&(\frac{1}{2}v^\top \beta_i\alpha v)_{k+1\leq i\leq n}\\
    0&0&g&\alpha v\\
    0&0&0&1
\end{pmatrix}
 \right\},
\end{align*}
where $(g,\mu(g))\in G(R),v\in \M_{2d,1}(R), z\in\M_{k,1}(R)$ for a certain subgroup $G$ of $\GSp(\Psi)$. Note the analogy of the group $Q$ with the group $Q_{\Lambda}$ in Lemma \ref{MT}. Up to conjugation, this will be the typical form of the algebraic group appearing in the mixed Shimura datum of a special subvariety of $\pgn$. The lower right $(n-k+2g+1)\times (n-k+2g+1)$ submatrix of it will account for the “Ribet section” part of the special subvariety.
\end{remark}
 
\section{The Tales of Zilber-Pink and Relative Manin-Mumford}\label{section6}

\subsection{The Zilber--Pink Conjecture for mixed Shimura varieties}

Let $S$ be a mixed Shimura variety over $\CC$ (may have several connected components). A \textbf{special subvariety} of $S$ is an irreducible component of a mixed Shimura subvariety of $S$, or of its image under a Hecke operator. (Hecke operators will transfer a subvariety of $S$ between different connected components of $S$. In previous sections of this article, all mixed Shimura varieties we considered are so-called \textit{connected} mixed Shimura varieties. So a special subvariety is just an irreducible component of a mixed Shimura subvariety, without the Hecke operator issues.) 

The Zilber-Pink Conjecture for mixed Shimura varieties is a statement about unlikely intersections in any mixed Shimura variety $S$. It concerns intersections of any irreducible closed subvariety $Z$ of $S$ with special subvarieties of $S$. 
implies the Mordell-Lang conjecture, the Andr\'{e}-Oort conjecture, and the Andr\'{e}-Pink conjecture \cite{pink04}. Three equivalent formulations of the Zilber-Pink conjecture can be found in \cite{Pink05}: \cite[Conjecture 1.1]{Pink05}, \cite[Conjecture 1.2]{Pink05}, and \cite[Conjecture 1.3]{Pink05}. We recall the formulation  \cite[Conjecture 1.3]{Pink05}. Recall that $Z$ is called a  \textbf{Hodge generic} subvariety of $S$, if  $Z$ is not contained in any proper special subvariety of $S$.

\begin{conjecture}[{\cite[Conjecture 1.3]{Pink05}}]\label{conj1.4}
Consider a mixed Shimura variety $S$ over $\CC$, and a Hodge generic irreducible closed subvariety $Z$. The union 
\[\bigcup_{\substack{\text{special subvariety } T\subset S\\ \operatorname{codim}_S T>\dim Z}}T\cap Z\]
is not Zariski dense in $Z$.
\end{conjecture}

One can formulate the analogous conjecture for semi-abelian varieties, where special subvarieties are precisely the irreducible components of subgroups. Then one gets the Zilber-Pink conjecture for semi-abelian varieties \cite[Conjecture 5.1]{Pink05}, which we will recall as Conjecture~\ref{Pink-Manin-Mumford}. Furthermore, Pink proposed a relative version of \cite[Conjecture 5.1]{Pink05}, which is \cite[Conjecture 6.1]{Pink05}.

It is natural to expect for Zilber-Pink for mixed Shimura varieties implies Zilber-Pink for semi-abelian varieties and its relative version, \textit{i.e.} Conjecture~\ref{conj1.4} implies both \cite[Conjecture 5.1]{Pink05} and \cite[Conjecture 6.1]{Pink05}. Pink claimed this in \cite{Pink05}. 
 However, as indicated in \cite{BE20}, Pink's arguments \cite[Theorem 5.7, Theorem 6.3]{Pink05} giving these deductions were incorrect. Further, the statement of \cite[Conjecture 6.1]{Pink05} was wrong, as the arguments in \cite{BE20} showed that the Ribet sections give counterexamples.

 

In the following subsections, we will make a correction based on our characterization of special subvarieties of $\pgn$  given by Theorem~\ref{general}. More preciesly, we will show that Conjecture~\ref{conj1.4} implies \cite[Conjecture 5.1]{Pink05}, and give the correct statement for \cite[Conjecture 6.1]{Pink05} and the Relative Manin-Mumford Conjecture. 


%There is a diagram that shows the implications of the conjectures involved. 

%\begin{center}

%\begin{tikzcd}[row sep=10ex, column sep=2ex]
%                                                      &  & \text{Pink Conjectures} \arrow[d, "\text{Theorem \ref{pink-PML}}"description,Rightarrow] \arrow[rrd,"\text{\cite[Theorem 3.3]{Pink05}}"description, Rightarrow] \arrow[lld,"\text{Proposition \ref{pink-re_RMM}}"description,Rightarrow] &  &                                                       \\
%\text{Relative Manin-Mumford} \arrow[rrd, Rightarrow] &  & \text{Pink-Mordell-Lang} \arrow[d,"\text{Fact \ref{PMM-PML}}"description, Leftrightarrow] \arrow[rrd, leftrightarrow, "\text{analogue}"description, dotted]                &  & {\cite[\text{Conjecture 1.6}]{pink04}} \arrow[d, "\text{\cite[Theorem 5.4]{pink04}}"description,Rightarrow] \\
%                                                      &  & \text{Pink-Manin-Mumford}                                                                     &  & \text{Mordell-Lang}                                  
%\end{tikzcd}

%\end{center}

\subsection{Zilber-Pink for mixed Shimura varieties implies Zilber-Pink for semi-abelian varieties}\label{SubsectionZPimpliesZPSA}
We first introduce the following notations.
\begin{notation}
    For an commutative algebraic group $G$ over $\CC$, we set $G^{[> d]}$ to be the union of all algebraic subgroups of $G$ of codimension $> d$, that is: \[
    G^{[> d]}:=\bigcup_{\substack{\text{algebraic subgroup } H\subset G\\ \operatorname{codim}_G H>d}}H.
    \]
    More generally, for an algebraic family of commutative algebraic groups $\calG \to S$, where $S$ is a variety over $\CC$, we set $\calG^{[> d]}:=\cup_{\substack{s\in S}}\calG_s^{[>d]}$.
\end{notation}

Pink discussed the following analog of the Manin-Mumford Conjecture and Mordell-Lang conjecture in \cite{Pink05}. 

\begin{conjecture}[Zilber-Pink for semi-abelian varieties {\cite[Conjecture 5.1]{Pink05}}]\label{Pink-Manin-Mumford} Consider a semi-abelian variety $G$ over $\CC$, and an irreducible closed subvariety $Z$, with $\dim Z=d$. Assume $Z$ is \textbf{not contained in any proper algebraic subgroup} of $G$. Then $Z\cap G^{[>d]}$ is not Zariski dense in $Z$.
\end{conjecture}
By \cite[Theorem 5.3, Theorem 5.5]{Pink05}, Conjecture \ref{Pink-Manin-Mumford} is equivalent to:

\begin{conjecture}[{\cite[Conjecture 5.2]{Pink05}}]\label{Pink-Mordell-Lang}
Consider a semi-abelian variety $G$ over $\CC$, a subgroup $\Gamma\subset G(\CC)$ of finite $\ZZ$-rank, and an irreducible closed subvariety $Z$, with $\dim Z=d$. Assume $Z$ is \textbf{not contained in any translation of proper algebraic subgroup} of $G$. Then $Z\cap (\Gamma+G^{[>d]})$ is not Zariski dense in $Z$.
\end{conjecture}

\begin{remark}
%    These conjectures are only \textit{analogs} of the original ones. In other words, the original conjectures are not direct consequences of the “generalized” conjectures proposed by Pink. The readers should be careful that the condition “not contained in any proper algebraic subgroup” is stronger than “not any algebraic subgroup”. Moreover, the conclusions are stronger than the original ones, since 
For closed proper subvariety $Z\subset G$ one has $d<\dim G$,\[
    \left\{ \text{torsion points}\right\}=G^{[>\dim G-1]} \subset G^{[>d]}.
    \]
So the Zilber-Pink Conjecture for semi-abelian varieties implies the 
%    we will see that, in the case of a family of semi-abelian varieties $\calG\to S$, the difference will lead to the conclusion “$\dim Z>\dim_S G$” of the Relative 
Manin-Mumford conjecture.
\end{remark}

To connect these conjectures with Pink's conjectures, we need to switch from the statement about a single semi-abelian variety to a moduli theoretic statement. Consider the semi-abelian variety $G$ in a family $\calG\to S$, with $\calG_s=G$. A typical such family is called a \textbf{Shimura family of semi-abelian varieties}.

\begin{definition}\label{Shi_family_of_semiav}
    Let $S_g$ be the moduli space of abelian varieties with a fixed polarization type and a sufficiently high level structure. A Shimura family of semi-abelian varieties is the base change of the universal semi-abelian scheme $\pgn\rightarrow (\Agv)^{[n]}$ to a special subvariety $S$ of $(\Agv)^{[n]}$.
\end{definition}

\begin{remark}
    Every semi-abelian variety is parametrized in a “minimal” Shimura family. Explicitly, for a semi-abelian variety $G$ over $\CC$, there is a Shimura family of semi-abelian varieties $\pi:\calG\to S$, with a Hodge generic point $s\in S(\CC)$, such that $G\cong \calG_s$. In fact, the existence of such a family (without assuming $s$ is Hodge generic) is clear. The only thing to do is to replace the base by the special closure of $s$.
\end{remark}

The advantage of viewing a semi-abelian variety $G$ in a Shimura family is that we can talk about \textbf{special subvariety} of $G$.

\begin{definition}
    Let $\calG\to S$ be a Shimura family of semi-abelian varieties with an identification $\calG_s= G$ as above, where $s\in S(\CC)$ is Hodge generic. An irreducible closed subvariety $T\hookrightarrow G$ is called \textbf{special} if $T$ is an irreducible component of $T'\cap\calG_s$ for some special subvariety $T'$ of $\calG$.
\end{definition}

We now give a correction of \cite[Theorem 5.7]{Pink05}. The original proof has an error in the sentence “Conversely, for any special subvariety $T\subset A$, every irreducible component of $T\cap A_s$ is a translate of a semi-abelian subvariety of $A_s$ by a torsion point.” Here $A$ is a Shimura family of semi-abelian varieties. The appearance of Ribet sections disproves this claim. However, this claim is not far from the truth, as the following proposition shows.

\begin{proposition}\label{fiber=subgp}
Let $T\hookrightarrow\pgn$ be a special subvariety. Then for each fiber over $s\in (\Agv)^{[n]}(\CC)$, $T_s$ is a finite union of translates of algebraic subgroups of $(\pgn)_s$.
\end{proposition}
\begin{proof}
First, we reduce to the case when $n=1$ and $T$ is a Ribet section. 

\textbf{(1) Reduction steps.}

Recall the result in Theorem \ref{general}. There exist special subvarieties $\{T_i\}_{1\leq i\leq n}$ of $\Pt$ whose images in $\Ag$ are the same, and $T$ is an irreducible component of the inverse image of $T_1\times_{\Ag}T_2\times_{\Ag}\cdots\times_{\Ag}T_n$ under a matrix isogeny $[h]$ of $\pgn$ over $\Ag$, where $h\in\GL_{n}(\QQ)\cap \M_{n,n}(\ZZ)$.

\underline{1. Clear the isogeny.} One can describe the effect of a matrix isogeny $[h]$ of $\pgn$ over $\Ag$ on a fiber $(\pgn)_s$ for $s\in (\Agv)^{[n]}(\CC)$. Denote the image of $s$ in $\ag$ by $x$. Then $s$ determines an extension:
\[\begin{tikzcd}
0 \arrow[r] & \gm^n \arrow[r] & (\pgn)_s \arrow[r] & {(\Ag)_{x}} \arrow[r] & 0
\end{tikzcd}.\]
The isogeny $[h]$ maps $(\pgn)_s$ into $(\pgn)_{[h](s)}$, which is the middle vertical arrow in the following pushout diagram.
\[\begin{tikzcd}
0 \arrow[r] & \gm^n \arrow[r] \arrow[d, "{[h]}"'] \arrow[dr, phantom, "\lrcorner", very near end] & (\pgn)_s \arrow[r] \arrow[d,d, "{[h]}"'] & ({\Ag})_{x} \arrow[r] \arrow[d] & 0  \\
0 \arrow[r] & \gm^n \arrow[r]                     & {(\pgn)_{[h](s)}} \arrow[r]  & (\Ag)_x \arrow[r]               & 0
\end{tikzcd}\]
Since this arrow is a homomorphism, $T_s$ is a finite union of translates of algebraic subgroups of $(\pgn)_s$ if and only if $T_{[h](s)}$ is a finite union of translates of algebraic subgroups of $(\pgn)_{[h](s)}$. Hence one can perform the isogeny $[h]$ and assume that $T$ is of the form $T_1\times_{\Ag}T_2\times_{\Ag}\cdots\times_{\Ag}T_n$.

\underline{2. Reduce to the $n=1$ case.} Recall the construction of $\pgn$. On each fiber $(\pgn)_{s}$, it is defined by the following pullback diagram.
\[\begin{tikzcd}
0 \arrow[r] & \gm^n \arrow[r] \arrow[d] & (\pgn)_s \arrow[r] \arrow[d, hook] \arrow[dr, phantom,"\ulcorner",very near start] & ({\Ag})_{x} \arrow[r] \arrow[d, "\Delta"', hook] & 0  \\
0 \arrow[r] & \gm^n \arrow[r]                     & (\Pt)^n_{s} \arrow[r]              & (\Ag)^n_x \arrow[r]                              & 0
\end{tikzcd}\]The middle vertical arrow is again a homomorphism. Hence to see if $T_s$ is a finite union of translates of algebraic subgroups of $(\pgn)_s$ is equivalent to see if its image in $(\Pt)^n_{s}$ is so. Since the latter is a product, it suffices to treat each factor, and this immediately reduces the question to the $n=1$ case.

\underline{3. Clear the torsion translation.} Now assume $T$ is a special subvariety of $\Pt$ and $s\in \Agv(\CC)$ whose image in $\ag$ is $x$. Recall that one can do the torsion translation process using the group laws of $\Pt$ over $\Ag$ and $\Agv$ respectively. The translation over $\Agv$ is easy to deal with, since the object $T_s$ is in $(\Pt)_s$, and translation over $\Agv$ preserves this fiber. The translation over $\Ag$ is slightly trickier since it no longer preserves the fiber $(\Pt)_s$. However, one can still describe how $(\Pt)_s$ transforms using certain Baer sum construction.

Now assume $y\in \Agv(\CC)$ is a torsion point whose image in $\ag$ is $x$. Recall that the extension
\[\begin{tikzcd}
0 \arrow[r] & \gm \arrow[r] & (\Pt)_y \arrow[r] & ({\Ag})_{x} \arrow[r] & 0
\end{tikzcd}\]
splits after a pullback along $(\Ag)_{x}$ by multiplying a large integer $N$, as indicated in the following diagram.
\[\begin{tikzcd}
0 \arrow[r] & \gm \arrow[d] \arrow[r] & {[N]^{*}(\Pt)_y} \arrow[d] \arrow[r] \arrow[dr, phantom,"\ulcorner",very near start] & (\Ag)_x \arrow[d, "{[N]}"'] \arrow[r] \arrow[l, "t"', dashed, bend right] & 0 \\
0 \arrow[r] & \gm \arrow[r]           & (\Pt)_y \arrow[r]                    & ({\Ag})_{x} \arrow[r]                                                     & 0
\end{tikzcd}\]
Denote the image of $t$ by $\tilde{t}$, which is an algebraic subgroup. Consider the following diagram.
\[\begin{tikzcd}
0 \arrow[r] & \gm^2 \arrow[d] \arrow[r]       & {(\Pt)_s\times [N]^{*}(\Pt)_y} \arrow[d] \arrow[r] \arrow[r] \arrow[dr, phantom,"\ulcorner",very near start] & (\Ag)_x^2 \arrow[d, "{[N]}"'] \arrow[r]     & 0 \\
0 \arrow[r] & \gm^2 \arrow[r] \arrow[d, "m"'] \arrow[dr, phantom, "\lrcorner", very near end] & (\Pt)_s\times(\Pt)_y \arrow[r] \arrow[d]           & ({\Ag})_{x}^2 \arrow[r] \arrow[d]           & 0 \\
0 \arrow[r] & \gm \arrow[r]                   & \widetilde{(\Pt)_{s+y}} \arrow[r]                        & (\Ag)_x^2 \arrow[r] \arrow[dl, phantom, "\llcorner", very near end]                         & 0 \\
0 \arrow[r] & \gm \arrow[u] \arrow[r]         & (\Pt)_{s+y} \arrow[u, hook] \arrow[r]                & (\Ag)_x \arrow[u, "\Delta", hook] \arrow[r] & 0
\end{tikzcd}\]
Here the last three arrows express that $(\Pt)_{s+y}$ is the Baer sum of $(\Pt)_{s}$ and $(\Pt)_{y}$. Consider the image of $[N]^*T_s\times \tilde{t}$ along the composition of the first two middle vertical arrows. It lies in $(\Pt)_{s+t}$, which is just the translation of $T_s$ by $\tilde{t}$. Since the first two middle vertical arrows in the above diagram are group homomorphisms, this translation process does not affect the property we want to show. Hence after performing suitable translations, we can assume that the image of $T$ in $\Ag\times_{\ag}\Agv$ is an abelian subscheme.

\underline{4. The case of sub-$\gm$-torsors.} We now deal with this easy case. Let $\calA'\subset \Ag\times_{\ag}\Agv$ be the image of $T$ in $\Ag\times_{\ag}\Agv$, which is an abelian subscheme (possibly with a smaller base). Now assume $T=\pi^{-1}(\calA')$, where $\pi: \Pt\rightarrow \Ag\times_{\ag}\Agv$ is the projection. Then for $s\in\Agv(\CC)$ whose image in $\ag$ is $x$, $((\Ag)_x \times \{s\})\cap \calA'_x$ is a tranlation of an abelian subvariety of $(\Ag)_x\times \{s\}$. Now it is easy to see that $T_s$ is a translation of a semi-abelian subvariety of $(\Pt)_s$.

\textbf{(2) The case of Ribet sections.}

Now assume $\calA'\subset \Ag\times_{\ag}\Agv$ is an abelian subscheme satisfying the Ribet property, and $T$ is the Ribet section over it. Recall that $T_x$ is constructed by the following diagram in Section \ref{section2}.
\[
\begin{tikzcd}
{[2]_{\calA_x'}^*((\Pt)_x|_{\calA'_x})} \arrow[d] \arrow[r]\arrow[dr, phantom,"\ulcorner",very near start]            & (\Pt)_x|_{\calA'_x} \arrow[d] \arrow[r]\arrow[dr, phantom,"\ulcorner",very near start] & (\Pt)_x \arrow[d]    \\
\calA'_x \arrow[r, "{[2]_{\calA'_x}}"] \arrow[u, "\tilde{r}_{\calA'_x}", dotted, bend left] & \calA_x' \arrow[r, hook]              & (\Ag)_x\times (\Agv)_x
\end{tikzcd}
\] To show that $T_s$ is a finite union of translates of algebraic subgroups, it suffices to show that for each $\tilde{s}\in [2]_{(\Agv)_x}^{-1}(s)$, the image $\tilde{r}_{\calA'_x} (((\Ag)_x\times \tilde{s})\cap \calA'_x)$ is so. Let $K\subset \calA'_x$ be the kernel of the projection $\calA'_x\rightarrow(\Agv)_x$. Then $\tilde{r}_{\calA'_x} (((\Ag)_x\times \tilde{s})\cap \calA'_x)$ is a finite union of translations of a section of $[2]_{\calA_x'}^*((\Pt)_x|_{\calA'_x})$ over $K\times \tilde{s}$. We claim that this section over $K\times \tilde{s}$ is a translation of an algebraic subgroup of $(\Pt)_s$. The proposition follows once this claim is proved.

By definition, the $\gm$-torsor $[2]_{\calA_x'}^*((\Pt)_x|_{\calA'_x})$ is trivial over $\calA'_x$ and in particular is trivial over $((\Ag)_x\times \tilde{s})\cap \calA'_x$. The claim follows if one can show that $[2]_{\calA_x'}^*((\Pt)_x|_{\calA'_x})$ is also trivial over $K\times \tilde{s}$, since then every section of $[2]_{\calA_x'}^*((\Pt)_x|_{\calA'_x})$ over $K\times\tilde{s}$ is a translation of an algebraic subgroup. Let $\lb_{\tilde{s}}\in \operatorname{Pic}^0((\Ag)_x\times\tilde{s})$ be the line bundle on $(\Ag)_x\times \tilde{s}$ represented by $\tilde{s}$. Suppose $((\Ag)_x\times \tilde{s})\cap \calA'_x=\bigcup_{i=1}^m (K+t_i)\times\tilde{s}$, then $\lb_{\tilde{s}}|_{(K+t_i)\times \tilde{s}}$ is trivial for $1\leq i\leq n$. Note that since $\lb_{\tilde{s}}\in \operatorname{Pic}^0((\Ag)_x\times\tilde{s})$, we have $t_{-t_1}^{*}\lb_{\tilde{s}}\cong \lb_{\tilde{s}}$. Hence
\[\lb_{\tilde{s}}|_{K\times \tilde{s}}\cong (t_{-t_1}^{*}\lb_{\tilde{s}})|_{(K+t_1)\times \tilde{s}}\cong \lb_{\tilde{s}}|_{(K+t_1)\times \tilde{s}}\]
is trivial and the claim follows.
\end{proof}

%Now we are going to prove that Pink's conjectures imply the Pink-Mordell-Lang conjecture. The original treatment in \cite{Pink05} is to deduce the Pink-Manin-Mumford conjecture from Pinks conjectures, and then use Fact \ref{1.1equiv1.2} to deduce the Pink-Mordell-Lang conjecture. Here we choose to directly show that Pink's conjectures imply the Pink-Mordell-Lang conjecture, the possibility of which has already been proposed in \cite[Remark 1.3]{BE20}.

\begin{theorem}[Zilber-Pink for mixed Shimura varieties implies Zilber-Pink for semi-abelian varieties]\label{pink-PML} 
Conjecture \ref{conj1.4} implies Conjeture \ref{Pink-Mordell-Lang}, and hence Conjecture~\ref{Pink-Manin-Mumford}. 
\end{theorem}

\begin{proof}

    %\begin{lemma}
       % Assume Conjecture \ref{conj1.4}. Then for any semi-abelian variety $G$ over $\CC$. if an irreducible closed subvariety $Z\hookrightarrow G$ satisfy the condition of Pink-Mordell-Lang, then the conclusion of Pink-Manin-Mumford holds for $Z$. 
        
       % More precisely, If $Z$ is not contained in any translation of algebraic subgroup of $G$, then $Z\cap G^{[>d]}$ is not Zariski dense in $Z$.
   % \end{lemma}
    
    %\begin{proof}
       % Take a Shimura family of semi-abelian varieties $\pi:\calG\to S$, with a Hodge generic point $s\in S(\CC)$, such that $G\cong \calG_s$. Under the identification, we regard $Z\hookrightarrow\calG_s\hookrightarrow\calG$ as an irreducible closed subvariety of $\calG$.
        
      %  By Proposition \ref{fiber=subgp}, any special subvariety $T$ of $\calG$, $T \cap \calG_s$ is a finite union of torsion translate of algebraic subgroups of $\calG_s$. Since $Z$ is not contained in any translation of an algebraic subgroup of $\calG_s$ as in the setting of Conjecture \ref{Pink-Mordell-Lang}, one immediately concludes that $Z$ is Hodge generic in $\calG$. 

      %  Meanwhile, since $s$ is Hodge generic in $S$, every semi-abelian subvariety of $\calG_s$ extends to a Shimura family of semi-abelian subvarieties over a finite covering of $S$. \textbf{Reference to be found} Also, for a $N$-torsion point $t\in \calG_s(\CC)$, a finite base change by $\calG[N]\to S$ would extend $t$ to a Shimura morphism $S\to \calG$ over $S$. Since any algebraic subgroup of $\calG_s$ is a finite union of torsion translation of semi-abelian varieties of $\calG_s$, we have thus an inclusion:
        %    \[\bigcup_{\substack{\text{algebraic subgroup } T\subset \calG_s\\ \operatorname{codim}_{\calG_s} T>\dim Z}}(T\cap Z)\subset\bigcup_{\substack{\text{special subvariety } T\subset \calG\\ \operatorname{codim}_\calG T>\dim Z}}(T\cap Z).\]
      %  Apply Pink Conjecture \ref{conj1.4} to the mixed Shimura variety $\calG$, the latter is not Zariski dense in $Z$, and so does the former. That is, we achieve the conclusion of the Pink-Manin-Mumford Conjecture \ref{Pink-Manin-Mumford}. 
  %  \end{proof}
     
     We will adapt the trick in \cite[Theorem 5.3]{Pink05} to apply to our situation. Let $\Gamma\subset G(\CC)$, $Z\hookrightarrow G$, $d=\dim Z$ be as in the statement of the Pink-Mordell-Lang Conjecture \ref{Pink-Mordell-Lang}. Fix a maximal sequence of linearly independent elements $\{a_1,\cdots, a_r\}\subset\Gamma$. 

        Let $C$ denote the Zariski closure of the subgroup of
    $G^r$ generated by the point $\underline a := (a_1, \cdots, a_r)$. After multiplying all $a_i$ by a suitable positive integer, if necessary, we may assume that $C$ is connected and hence a semi-abelian subvariety of $G^r$. Consider the semi-abelian variety $H := G \times C$ and its irreducible closed subvariety $Y := Z \times  \{\underline a\}$ of dimension $d$. 
    
    Different from the original proof, we will show that $Y$ is not contained in any \textit{proper special subvariety} of $H$. Let $n$ be the toric rank of $H$, $g=\dim H-n$. Then there is $s\in (\Agv)^{[n]}(\CC)$ such that $H\cong(\pgn)_{s}$. Let $x\in \ag$ be the image of $s$ under the natural projection. Assume there is a special subvariety $T$ of $\pgn$ such that $Y\subset T_s=T\cap (\pgn)_{s}$. 

    Since $Z$ is not contained in any translation of any algebraic subgroup of $G$, the differences of elements in $Z$ generate $H$. The smallest algebraic subgroup of $H$ containing $Y$ thus contains $G\times \{\underline{a}\}$. Since $\{\underline{a}\}$ is Zariski dense in $C$, this subgroup of $H$ must be the whole $H$. In other words, $Y$ is not contained in any proper algebraic subgroup of $H$.

    The above arguments imply that the special closure of the image of $Y$ in $(\Ag)_x$ is the whole $(\Ag)_x$. Thus the image of $T_s$ in $(\Ag)_x\times (\Agv)^{[n]}_x$ contains $(\Ag)_x\times \{s\}$. It is also a torsion translation of an abelian subvariety $A'$ of $(\Ag)_x\times (\Agv)^{[n]}_x$. 
    
    We claim that at this stage, $T_s$ must be a finite union of torsion translations of algebraic subgroups of $H$, hence must be the whole $H$. Recall the classifications in Theorem \ref{general}. The key is to show that non-trivial Ribet sections do not appear. Assume that up to a matrix isogeny of $(\Agv)^{[n]}_x$, there exists $1\leq i\leq n$ such that $(\operatorname{id}_{(\Ag)_x}\times \operatorname{pr}_i)(A')$ is an abelian subvariety of $(\Ag)_x\times (\Agv)_x$ satisfying the Ribet property. However, $(\operatorname{id}_{(\Ag)_x}\times \operatorname{pr}_i)(A')$ also contains a subvariety of $(\Ag)_x\times (\Agv)_x$ of the form $(\Ag)_x\times \tilde{s}$. Corollary \ref{not_entire} then shows that $(\operatorname{id}_{(\Ag)_x}\times \operatorname{pr}_i)(A')$ must be $(\Ag)_x\times 0$.

    Now apply Pink conjectures to $Y\subset \pgn$, we now deduce that $Y\cap H^{[>d]}$ is not Zariski dense in $Y$. The arguments in \cite[Theorem 5.3]{Pink05} showed that 
    \[Z\cap (\Gamma+G^{[>d]})\times \{\underline{a}\}\subset Y\cap H^{[>d]}.\] Hence $Z\cap (\Gamma+G^{[>d]})$ is not Zariski dense in $Z$, which is the conclusion in the Pink-Mordell-Lang conjecture.
\end{proof}

\subsection{Reformulation of the Relative Manin-Mumford Conjecture}\label{SubsectionReformulationRMM}

In this section, we assume that $S$ is a variety over $\CC$ throughout.

Pink made the following relative version of Conjecture \ref{Pink-Manin-Mumford}. Its original statement is as follows.
\begin{conjecture}[{\cite[Conjecture 6.1]{Pink05}}]\label{RMM}
Consider an algebraic family of semi-abelian varieties $\calG\rightarrow S$ over $\CC$ and an irreducible closed subvariety $Z\subset \calG$ of dimension $d$ that is not contained in any proper closed subgroup scheme (with a possibly smaller base) of $\calG\rightarrow S$. Then $Z\cap \calG^{[>d]}$ is not Zariski dense in $Z$.
\end{conjecture}
Now suppose the base $S$ is irreducible, then the family $\calG\rightarrow S$ has a constant relative dimension, say $n$. Then for any $d<n$, the subset $\calG^{[>d]}$ contains all torsion points of all fibers of $\calG\rightarrow S$. Hence the following is an easy consequence of Conjecture \ref{RMM}.
\begin{conjecture}[{\cite[Conjecture 6.2]{Pink05}}]\label{RMMcor}

    Consider an algebraic family of semi-abelian varieties $\calG\rightarrow S$ of relative dimension $n$ over an irreducible variety over $\CC$, and an irreducible closed subvariety $Z\subset \calG$. Assume that $Z$ is not contained in any proper closed subgroup scheme of $\calG\rightarrow S$ and that it contains a Zariski dense subset of torsion points. Then $\dim Z\geq n$.

\end{conjecture}

Conjecture~\ref{RMMcor} is called the {\em Relative Manin-Mumford Conjecture}. 
When the family $\calG\rightarrow S$ is a family of abelian varieties, Conjecture \ref{RMMcor} is recently proved in \cite{GaoHab_RMM}. However, when the family $\calG\rightarrow S$ is a family of semi-abelian varieties of non-zero toric rank, a counterexample to Conjecture \ref{RMM} is constructed in \cite{BE20} (see the discussions after \cite[Theorem 2.4]{BE20}). 

Briefly, the family they considered is the Poincar\'{e} torsor over $E$ for a CM elliptic curve $E$ over $\CC$. Let $\lambda:E\stackrel{\sim}{\rightarrow} E^{\vee}$ be the canonical principal polarization on $E$ given by the line bundle $\mathcal{O}_E(0_E)$. Choose an endomorphism $\phi\in\operatorname{End}_\CC(E)$ with $\phi^{\dagger}\neq \phi$, and set $\alpha:=\lambda\circ(\phi-\phi^{\dagger})$. They constructed the Ribet section $t_{\phi}$ over the graph of $\alpha$ and showed that 
\begin{itemize}
    \item the image of $t_{\phi}$ is not contained in a proper subgroup scheme of $\mathcal{P}_E^{\times}$,
    \item and $t_{\phi}(x,\alpha(x))$ is torsion in the fiber of $\mathcal{P}_E^{\times}\rightarrow E$ for every torsion point $x$ of $E$.
\end{itemize} Now let $Z=t_{\phi}(\Gamma(\alpha))$ and $d=1$ in Conjecture \ref{RMM}. In this case, $\mathcal{P}_E^{\times,[>1]}$ is the set of points of $\mathcal{P}_E^{\times}$ that are torsion in their fiber and $Z\cap\mathcal{P}_E^{\times,[>1]}$ is dense in $Z$, hence it is a counterexample. On the other hand, $t_{\phi}$ is still a special subvariety of $\mathcal{P}_E^{\times}$ as a mixed Shimura variety.

This counterexample indicates that the formulation of the Relative Manin-Mumford conjecture in the semi-abelian scheme case is incorrect, and a problem is that there are more “special subvarieties” of a family of semi-abelian varieties than just torsion translates of subgroup schemes. On the other hand, our classification of the special subvarieties of the universal semi-abelian scheme showed that they are the only exceptions. Hence one would include the Ribet sections in the statement of the Relative Manin-Mumford conjecture and formulate a  more plausible version.% of the Relative Manin-Mumford conjecture.

To achieve this goal, we define first what a “special subvariety” of a family of semi-abelian varieties is. Let $\calG\to S$ be a semi-abelian scheme of relative dimension $g+n$, with toric rank $n$, and $Z\subset \calG$ be an irreducible closed subvariety. we are going to define the specialty of $Z$.

\begin{definition}\label{moddef}
    Let $\calG/S$ be a semi-abelian scheme of relative dimension $g+n$, with toric rank $n$. Let $Z$ be an irreducible closed subvariety of $\calG$. Suppose $Z$ is an irreducible component of the pullback of a special subvariety of $\pgn$ via the classifying morphism $S\to (\Agv)^{[n]}$ (up to a finite \'{e}tale base change to gain a sufficiently high level structure). In this case, we say $Z$ is \textbf{special}.
\end{definition}

\begin{remark}
Note that the above definition is “moduli-theoretic” in essence. It is possible to translate it into a purely intrinsic definition, but this would be too elaborate. As a side remark, we invite the reader to translate the definitions in \cite{Known_for_semiab_surf} to our definition, and get a feeling of how elaborate an intrinsic definition could be.
\end{remark}

We can now formulate the conjectures, parallel to the original ones formulated by Pink. We start with the revised version of \cite[Conjecture 6.1]{Pink05}.
\begin{conjecture}\label{re_RMM}
    Let $\calG\rightarrow S$ be a semi-abelian scheme over $\CC$, $Z\subset \calG$ be an irreducible closed subvariety of dimension $d$ that is not contained in a proper special subvariety of $\calG$. Then $Z\cap \calG^{[>d]}$ is not Zariski dense in $Z$.
\end{conjecture}
Next we have the revised version of the Relative Manin-Mumford Conjecture. It is an immediate corollary of Conjecture~\ref{re_RMM}.

\begin{conjecture}[Relative Manin-Mumford, revised]\label{re_RMMcor}
    Let $\calG/S$ be a semi-abelian scheme of relative dimension $g+n$, with toric rank $n$. Let $Z$ be an irreducible closed subvariety of $\calG$, dominant over $S$. If $Z(\CC) \cap \calG_{\text{tor}}$ is Zariski dense in $X$, then either $Z$ is contained in a special subvariety of $\calG$, or $\dim X \geq g+n$.
\end{conjecture}

\begin{proposition}[Zilber--Pink imples (revised) Relative Manin-Mumford]\label{pink-re_RMM}
    Conjecture~\ref{conj1.4} implies Conjecture \ref{re_RMM}, and hence Conjecture~\ref{re_RMMcor}.
\end{proposition}
\begin{proof}
Note that Conjecture \ref{re_RMM} is invariant under a finite \'{e}tale base change of the base. Hence up to a finite \'{e}tale base change to equip the family with a sufficiently high level structure, we can assume that the family $\calG\to S$ is parametrized in the universal Shimura family $\pgn\to (\Agv)^{[n]}$ of semi-abelian varieties via the classification morphism $S\to (\Agv)^{[n]}$. Restrict the family to the special closure of the schematic image of the classification morphism, and let $Z'$ be the schematic image of $Z$ under $\calG\to \pgn$. Since $Z$ is not contained in a proper special subvariety of $\calG$, $Z'$ is not contained in any proper special subvariety of $\pgn$, which means $Z'$ is Hodge generic in $\calG$. Apply Pink's Conjecture \ref{conj1.4}, we get the set \[
\bigcup_{\substack{\text{special subvariety } T\subset \pgn\\ \operatorname{codim}_{\pgn} T>\dim Z'}}(T\cap Z')
\]
is not Zariski closed in $Z'$. As $Z$ is the fibered product $Z'\times_{(\Agv)^{[n]}}S$, the corresponding set in $\calG$ is not Zariski dense in $Z$:
\[
\bigcup_{\substack{\text{special subvariety } T\subset \pgn\\ \operatorname{codim}_{\pgn} T>\dim Z'}}((T\times_{(\Agv)^{[n]}} S)\cap Z).
\]
We claim that the above set contains all algebraic subgroups of $\calG_s$ with codimension $>d$ for any $s\in S(\CC)$. If the claim is achieved, one immediately has $Z\cap \calG^{[>d]}$ is contained in a set which is not Zariski dense in $Z$, and the conclusion of Conjecture \ref{re_RMM} follows.

Fix $s\in S(\CC)$. Let $s'$ be the image of $s$ under $S\to (\Agv)^{[n]}$, and $S'$ be the special closure of $s'$. Every semi-abelian subvariety of $\calG_s=(\pgn)_{s'}$ extends to a Shimura family of semi-abelian subvarieties over a finite covering of $S'$ (by varying the level structure).\footnote{For a proof of the fact, one consider the orbit of any Hodge generic point under the Mumford-Tate group, and take the image of the period map.} Also, for an $N$-torsion point $t\in \calG_s(\CC)$, a finite base change by $(\pgn)_{S'}\to S'$ would extend $t$ to a Shimura morphism $S'\to (\pgn)_{S'}$ over $S'$. These two extension property realizes semi-abelian subvarieties and torsion sections as the pullback of special subvarieties. As every algebraic subgroup of $\calG_s$ is a finite union of torsion translation of a semi-abelian subvariety of it, and the extension properties preserve the dimension relations, we conclude that:
\[
\calG_s^{[>d]}\subset\bigcup_{\substack{\text{special subvariety } T\subset \pgn\\ \operatorname{codim}_{\pgn} T>\dim Z'}}(T\times_{(\Agv)^{[n]}} S).
\]
which proves the claim.
\end{proof}

%So far there are few results on the reformulated relative Manin-Mumford conjecture, largely due to the lack of a complete classification of special subvarieties of the universal semi-abelian variety. However, previously there are some classification results in the case of the universal Poincar\'{e} torsor over the modular curve. So there are definite results in the case when the base is a curve and the semi-abelian varieties are extensions of elliptic curves by $\gm$. See the work \cite{Known_for_semiab_surf}.

\newcommand{\etalchar}[1]{$^{#1}$}
\providecommand{\bysame}{\leavevmode\hbox to3em{\hrulefill}\thinspace}
\providecommand{\MR}{\relax\ifhmode\unskip\space\fi MR }
% \MRhref is called by the amsart/book/proc definition of \MR.
\providecommand{\MRhref}[2]{%
  \href{http://www.ams.org/mathscinet-getitem?mr=#1}{#2}
}
\providecommand{\href}[2]{#2}
\begin{thebibliography}{BMPZ13}

\bibitem[BE20]{BE20}
Daniel Bertrand and Bas Edixhoven, \emph{Pink's conjecture on unlikely intersections and families of semi-abelian varieties}, J. \'Ec. polytech. Math. \textbf{7} (2020), 711--742. \MR{4086583}

\bibitem[BMPZ13]{Known_for_semiab_surf}
Daniel Bertrand, David Masser, Anand Pillay, and Umberto Zannier, \emph{Relative manin-mumford for semi-abelian surfaces}, 2013.

\bibitem[Del74]{hodge3}
Pierre Deligne, \emph{Th\'eorie de {Hodge} : {III}}, Publications Math\'ematiques de l'IH\'ES \textbf{44} (1974), 5--77 (fr). \MR{498552}

\bibitem[Edi21]{EdixhovenPrivate}
Bas Edixhoven, \emph{On implications between pink's unlikely intersection conjectures}, unpublished manuscript (2021).

\bibitem[FC90]{faltings-chai}
Gerd Faltings and Ching-Li Chai, \emph{Degeneration of abelian varieties}, Ergebnisse der Mathematik und ihrer Grenzgebiete (3) [Results in Mathematics and Related Areas (3)], vol.~22, Springer-Verlag, Berlin, 1990, With an appendix by David Mumford. \MR{1083353}

\bibitem[Ful98]{fulton}
William Fulton, \emph{Intersection theory}, second ed., Ergebnisse der Mathematik und ihrer Grenzgebiete. 3. Folge. A Series of Modern Surveys in Mathematics [Results in Mathematics and Related Areas. 3rd Series. A Series of Modern Surveys in Mathematics], vol.~2, Springer-Verlag, Berlin, 1998. \MR{1644323}

\bibitem[Gao16]{GaomixAO}
Ziyang Gao, \emph{About the mixed andr{\'e}–oort conjecture: Reduction to a lower bound for the pure case}, Comptes Rendus Mathematique \textbf{354} (2016), no.~7, 659--663.

\bibitem[Gao17a]{GaoAL}
Z.~Gao, \emph{Towards the {A}ndr\'{e}-{O}ort conjecture for mixed {S}himura varieties: the {A}x-{L}indemann-{W}eierstrass theorem and lower bounds for {G}alois orbits of special points}, J.Reine Angew. Math (Crelle) \textbf{732} (2017), 85--146.

\bibitem[Gao17b]{gaoziyang}
Ziyang Gao, \emph{A special point problem of {A}ndr\'e-{P}ink-{Z}annier in the universal family of {A}belian varieties}, Ann. Sc. Norm. Super. Pisa Cl. Sci. (5) \textbf{17} (2017), no.~1, 231--266. \MR{3676047}

\bibitem[GH23]{GaoHab_RMM}
Ziyang Gao and Philipp Habegger, \emph{The relative manin-mumford conjecture}, arXiv e-prints (2023), \href{https://arxiv.org/abs/2303.05045}{\texttt{arXiv:2303.05045}}.

\bibitem[Mum70]{Mum12}
David Mumford, \emph{Abelian varieties}, Tata Institute of Fundamental Research Studies in Mathematics, vol.~5, Tata Institute of Fundamental Research, Bombay; by Oxford University Press, London, 1970. \MR{282985}

\bibitem[MvdGE]{BGB_AV}
Ben Moonen, Gerard van~der Geer, and Bas Edixhoven, \emph{Abelian varieties}, \href{http://van-der-geer.nl/~gerard/AV.pdf}{\texttt{http://van-der-geer.nl/$\sim$gerard/AV.pdf}}.

\bibitem[Pin89]{PinkThesis}
R.~Pink, \emph{Arithmetical compactification of mixed {S}himura varieties}, Ph.D. thesis, Bonner Mathematische Schriften, 1989.

\bibitem[Pin05a]{pink04}
Richard Pink, \emph{A combination of the conjectures of {M}ordell-{L}ang and {A}ndr\'e-{O}ort}, Geometric methods in algebra and number theory, Progr. Math., vol. 235, Birkh\"auser Boston, Boston, MA, 2005, pp.~251--282. \MR{2166087}

\bibitem[Pin05b]{Pink05}
\bysame, \emph{A common generalization of the conjectures of andr{\'e}-oort, manin-mumford, and mordell-lang}, 2005, \href{https://people.math.ethz.ch/~pink/ftp/AOMMML.pdf}{\texttt{https://people.math.ethz.ch/$\sim$pink/ftp/AOMMML.pdf}}.

\bibitem[PST{\etalchar{+}}21]{AO}
Jonathan {Pila}, Ananth~N. {Shankar}, Jacob {Tsimerman}, H{\'e}l{\`e}ne {Esnault}, and Michael {Groechenig}, \emph{{Canonical Heights on Shimura Varieties and the Andr{\'e}-Oort Conjecture}}, arXiv e-prints (2021), \href{https://arxiv.org/abs/2109.08788}{\texttt{arXiv:2109.08788}}.

\bibitem[{Ray}83]{Ray83}
M.~{Raynaud}, \emph{{Courbes sur une vari{\'e}t{\'e} ab{\'e}lienne et points de torsion.}}, Inventiones Mathematicae \textbf{71} (1983), no.~1, 207--233.

\bibitem[Tsi18]{Tsimerman}
Jacob Tsimerman, \emph{The {Andr{\'e}}-{Oort} conjecture for {{\(\mathcal A_g\)}}}, Ann. Math. (2) \textbf{187} (2018), no.~2, 379--390 (English).

\bibitem[Zan12]{Zan12}
Umberto Zannier, \emph{Some problems of unlikely intersections in arithmetic and geometry}, Princeton University Press, Princeton, 2012.

\bibitem[Zha98]{Zhang1998}
Shou\-{wu} Zhang, \emph{Small points and arakelov theory.}, Documenta Mathematica (1998), 217--225.

\end{thebibliography}
